% =====================================================================
%  capitulos/cap01.tex — Capítulo 1. La gramática y su enseñanza
%  Autora: Mabel Giammatteo
%  Fuente: Capitulo_1_La_gramatica_y_su_ensenanza_normalizado.docx
% =====================================================================
\capitulo{La gramática y su enseñanza}{Mabel Giammatteo}

\section{Introducción}

La lengua es una entidad que ni siquiera necesita a sus hablantes para existir,
como prueban los numerosos idiomas muertos que han sobrevivido a sus pueblos.
Los individuos, en cambio, no podemos hacer prácticamente nada ---o nada
relevante para la vida en comunidad--- sin utilizarla. Pero la lengua no es un
instrumento externo con el que debamos familiarizarnos o que tengamos que
aprender a manejar, como si fuera un celular, un electrodoméstico o un
automóvil, sino que, como sostiene \textcite[18]{bosque2018}, se trata de «un
sistema complejo que los hablantes tenemos interiorizado y que constituye una
parte esencial de nuestra naturaleza». Conforma una estructura que, junto con
el vocabulario, se encuentra almacenada en nuestra mente y a diario la
empleamos en nuestros intercambios. Con solo abrir la boca\ldots, podríamos
decir.

Comenzamos este capítulo introductorio con un recorrido evaluativo sobre
distintos planteos que a lo largo de la historia se han hecho respecto de la
gramática y su utilidad para la enseñanza. Revisamos así conceptos centrales
como los de lengua y habla, competencia y actuación, sistema y uso, entre
muchos otros. Luego, introducimos la propuesta léxico-gramatical que se irá
desarrollando a lo largo de todo este libro y que tiene como objetivo avanzar
desde las nociones teóricas a la reflexión sobre el lenguaje de todos los días.
Para ello, además de describir y explicar las estructuras lingüísticas, nos
interesa conectarlas con el uso. Esperamos así poder colaborar en el
reconocimiento de los valores expresivos y comunicativos de las configuraciones
morfosintácticas y léxico-semánticas que se irán estudiando en este libro.

\section{Debates e interrogantes sobre la gramática y su enseñanza}

Si la lengua forma parte del conocimiento internalizado de los hablantes, que
diariamente la utilizamos para comunicarnos a través de distintos medios y
redes sociales, ¿para qué se necesita enseñarla en la escuela? ¿Para qué
insistir con paradigmas verbales y complicados análisis de oraciones?
Detengámonos en esto. Como señaló \textcite{saussure1945}, el lenguaje tiene
un doble costado. Por un lado, es lo que cada uno de nosotros «dice» en cada
una de las distintas circunstancias en que lo empleamos, lo que denominó
\termino{habla} y caracterizó como momentánea, individual y variable. Pero
también, por otro lado, el lenguaje tiene un costado social, la
\termino{lengua} propiamente dicha, que constituye un factor de herencia,
cultura y tradición. La lengua es, sin duda, un componente de nuestra
identidad, que nos atraviesa, interpela y define. Como facultad humana propia
de nuestra especie \parencite{chomsky1989}, el lenguaje es universal, pero la
lengua es propia de cada cultura, de cada pueblo y también de cada hablante que
la hace suya. Desde esta perspectiva, la educación lingüística en la escuela
cumple primariamente una función social: enseñar la lengua estándar a
poblaciones muchas veces heterogéneas y que se expresan mediante diferentes
variedades de uso. Por supuesto, esto no quiere decir que estas formas
particulares del habla de cada grupo o región, también conocidas como
lectos,\footnote{\pendiente{Falta el texto de esta nota en el original.}}
deban ser abandonadas o menospreciadas, ya que su dominio enriquece y
caracteriza al que lo posee. Pero la incorporación de una variedad estándar
facilita el acercamiento a los textos que se necesita comprender y producir en
el ámbito escolar. Asimismo, la enseñanza también debe hacer consciente ese
conocimiento tácito y subliminal, que naturalmente el hablante posee sobre su
lengua. Por tanto, la escuela tiene que ofrecer lo siguiente:

\begin{itemize}
  \item un \termino{conocimiento del uso}, que respete la variación, pero que
        también incorpore los aspectos formales necesarios para acceder a los
        textos con los que se debe interactuar en los estudios;
  \item un \termino{conocimiento reflexivo de la estructura y posibilidades},
        que pueda no solo colaborar con la mejora de las habilidades
        cognitivo-lingüísticas ---argumentar, fundamentar, comparar, explicar,
        entre otras---, sino también incentivar el desarrollo de la capacidad
        de abstracción y el juicio crítico, que se requieren tanto para los
        estudios como para la vida cotidiana.
\end{itemize}

\begin{recuadro}[Síntesis]
Toda actividad reflexiva sobre la lengua redunda en un mayor conocimiento de
las posibilidades expresivas que ofrece, lo cual, hacia el exterior, se
proyecta en mejores intercambios comunicativos y, hacia el interior, incide en
la capacidad de conceptualización y de establecer relaciones entre fenómenos
del mundo circundante.
\end{recuadro}

En esta perspectiva, la enseñanza de la lengua materna en la escuela debe
apelar a un conocimiento metacognitivo\footnote{El término
\textit{metacognición} \parencite{flavell1979} refiere a la capacidad de
reflexionar sobre los propios procesos cognitivos, es decir, tener control
consciente sobre las operaciones mentales que realizamos.} que forme hablantes
conscientes, que no solo reconozcan las estructuras de la lengua, sino que
también puedan identificar los valores comunicativos que transmiten y avanzar
desde la interpretación de los significados convencionales, que son los que
tenemos internalizados, a otros más novedosos o complejos, que exigen mayor
esfuerzo interpretativo, como los sentidos metafóricos o los de mayor nivel de
abstracción.

Para ilustrar la problemática que estamos tratando, podemos partir de una
historieta publicada hace ya un tiempo en el diario \textit{Clarín} de Buenos
Aires, en la que tiene lugar el siguiente diálogo entre una madre y su hija:

\figurapendiente{Donatela, por Raznovich (\textit{Clarín}, 2/11/12)}{fig:donatela}

Más allá de la ironía presente en la historieta, \forma{comunicarnos} aquí
refiere a uno de los usos más básicos del lenguaje: el simple \sig{entrar en
contacto} con el otro, al que Jakobson (1984)\pendiente{Jakobson 1984 no está
en la bibliografía} denominó \termino{función fática},\footnote{Las funciones
del lenguaje, según Jakobson (1984), son seis: referencial, centrada en la
transmisión de información; emotiva, que manifiesta la expresividad del emisor;
conativa, que apela al receptor del mensaje; metalingüística, que se emplea
cuando usamos el código lingüístico para hacer referencia al lenguaje mismo;
poética, que pone el acento en la conformación misma del mensaje; y fática,
centrada en el funcionamiento del canal a través del cual se produce el
intercambio.} la cual es muy útil en el inicio de los estudios de una lengua
extranjera, cuando básicamente lo que queremos es «conectarnos» con sus
hablantes mediante saludos e intercambios corteses. Sin embargo, estudiar la
lengua materna implica una actitud muy diferente, ya que el objetivo no es
aprender a hablarla, sino que lo que se busca es adquirir la práctica de
reflexionar sobre los matices significativos que sus diferentes construcciones
permiten transmitir e interpretar.

Volviendo a la inclusión de la gramática en la enseñanza, cuestión que, como
todos sabemos, viene siendo debatida desde hace ya largo tiempo, las distintas
posturas han oscilado entre la negación absoluta de su utilidad, por un lado, y
encendidas defensas de su contribución a la formación lingüística, por el otro.
Por eso, antes de avanzar, debemos especificar a qué nos referimos con
\forma{gramática}, ya que, según veremos a continuación, a lo largo de la
historia, no siempre se entendió el término de la misma manera.

\subsection{De la gramática tradicional al estructuralismo}

En el panorama de los estudios de la Antigüedad clásica y el Medioevo, la
gramática, junto con la dialéctica y la retórica, conformaba el
Trivium,\footnote{El Trivium y el Quadrivium ---formado por aritmética,
geometría, astronomía y música--- constituían el primer y segundo ciclo de las
llamadas artes liberales, surgidas en el siglo \textsc{vi} y representativas
del saber de la época.} destinado a comprender «los secretos de las palabras».
Tal como transmite \textcite[108]{barthes1993}, las gramáticas latinas de
Donato (hacia 350) y Prisciano (fines del s.~\textsc{v} y comienzos del
s.~\textsc{vi}) «representaron la ley absoluta, salvo cuando no coinciden con
la Vulgata:\footnote{Es la «versión latina de la Sagrada Escritura, [\ldots]
declarada auténtica por la Iglesia católica» (\textit{DLE}, en línea).} la
gramática no podía ser entonces más que normativa,\footnote{La gramática
normativa o prescriptiva es la que se ocupa de los usos correctos o incorrectos
según la norma aceptada socialmente en un momento determinado. Estos conceptos
se amplían en el capítulo~2.} puesto que se pensaba que las \enquote*{reglas}
de la locución habían sido inventadas por los gramáticos [\ldots]. Hasta el
siglo \textsc{xii}, la \textit{Grammatica} abarca la gramática y la poesía,
trata a la vez de la \enquote*{precisión} y de la \enquote*{imaginación}; de las
letras, las sílabas, la oración, el período,\footnote{En gramática,
\textit{período} hace referencia a la «estructura oracional, generalmente
bimembre, formada por una oración principal y una subordinada» (\textit{DLE},
en línea), sobre todo las construcciones condicionales y concesivas.} las
figuras\footnote{Las figuras retóricas son empleos no convencionales de las
palabras con fines expresivos.} y la métrica».\footnote{Estudio de la medida o
estructura de los versos y de sus distintas combinaciones.}

Las teorías sobre el lenguaje basadas en los estudios de los gramáticos griegos
y latinos constituyen lo que se conoce como gramática tradicional, cuya
influencia en la enseñanza se mantuvo vigente hasta la primera mitad del
siglo~\textsc{xx}. Su enfoque estaba centrado en la palabra y en sus
«accidentes» ---variaciones en la terminación para manifestar las categorías
morfológicas (género, número, caso, persona, tiempo, modo, aspecto,
voz)\footnote{El concepto de categoría morfológica se trata en el capítulo~4.
Para ampliar, v. \textcite{giammatteo2018}.}---. La sintaxis ocupaba un lugar
bastante marginal, de modo que, en las gramáticas y manuales escolares, las
nociones sintácticas solían colocarse al final y se orientaban sobre todo a
favorecer la buena formación de las oraciones y a evitar los llamados «vicios»
o errores de construcción.\footnote{Entre los considerados vicios del lenguaje
se encuentran el barbarismo (uso de extranjerismos), el solecismo (alteración
del orden sintáctico), el pleonasmo (agregado de palabras innecesarias), la
anfibología (uso de palabras con doble sentido), la cacofonía (combinación
desagradable de sonidos en un sintagma), entre otros.}

\begin{destacado}
Para la gramática tradicional, el interés fundamental radicaba en el «buen uso
de la lengua», por lo que las reglas y preceptos eran de tipo normativo y no
había una explicitación del funcionamiento de la lengua.
\end{destacado}

En el siglo \textsc{xix} la lingüística se encontraba influida por el afán
cientificista del positivismo,\footnote{Teoría filosófica formulada por
Auguste Comte, para la cual la única fuente del conocimiento es la experiencia,
por lo que el conocimiento debe ser observable, verificable y objetivo.} que
buscaba para las humanidades el rigor metodológico de las ciencias naturales,
por lo que el estudio del lenguaje comenzó a ser «cuidadosa y objetivamente
investigado y luego explicado en términos de hipótesis inductivas»
\parencite[22]{lyons1985}. Asimismo, la atracción que ejercieron los planteos
evolucionistas de Darwin hizo que la observación y el análisis de los datos se
enfocaran al estudio histórico y comparativo de las lenguas. Las
investigaciones de la época se basaban fundamentalmente en el análisis de las
distintas lenguas de la familia indoeuropea,\footnote{La mayoría de las lenguas
de Europa (el griego, el latín y sus derivadas las lenguas romances, las
lenguas germánicas, las eslavas, las célticas, etc.) y algunas de Asia (el
sánscrito y otras lenguas indoiranias, el persa, etc.) pertenecen a esta
familia.} que, por provenir de un tronco común y poseer textos escritos,
constituían un excelente campo de aplicación de la metodología comparativa. Los
estudios se hacían contrastando las unidades lingüísticas entre sí,
considerando su origen común y las características de su desarrollo histórico
en cada lengua, pero sin ir más allá, es decir que no se contemplaban sus
relaciones respecto del funcionamiento global de la lengua.

La idea de que en la lengua todos los elementos están interrelacionados, es
decir, que constituyen un \termino{sistema}, surge recién a principios del
siglo~\textsc{xx}, a partir de los planteos de \textcite{saussure1945},
considerado el padre de la lingüística moderna y el iniciador de la corriente
conocida como estructuralismo.\footnote{No obstante, Saussure no usó la palabra
\textit{estructura}, sino \textit{sistema}.} Los planteos de Saussure
impulsaron un cambio de paradigma en los estudios sobre el lenguaje, que
pasaron a ocuparse del conocimiento

\begin{itemize}
  \item no solo de las lenguas de prestigio, sino de \textit{todas las lenguas};
  \item no solo de las formas correctas (lingüística prescriptiva), sino de
        \textit{todas las formas} (lingüística descriptiva);
  \item no solo de la lengua escrita, sino primordialmente de \textit{la lengua
        hablada};
  \item no solo del acto individual, sino de \textit{la naturaleza social del
        lenguaje};
  \item no solo del desarrollo histórico de las lenguas (\textit{diacronía}),
        sino de \textit{su estado actual} (\textit{sincronía});
  \item no de la lengua como una nomenclatura, constituida por elementos
        aislados, sino como \textit{un sistema de signos formados por un
        significado} (contenido) \textit{y un significante} (soporte material
        de la expresión).\footnote{El soporte material debe ser perceptible,
        como el sonido en la lengua hablada, los gestos en la lengua de señas,
        o los signos gráficos en la escritura.}
\end{itemize}

Para Saussure, la tarea de la lingüística debía centrarse en los siguientes
puntos:

\begin{itemize}
  \item describir todas las lenguas y familias de lenguas;
  \item formular leyes generales;
  \item delimitarse y definirse a sí misma, es decir, constituirse como
        disciplina autónoma e independiente.
\end{itemize}

\begin{destacado}
Frente al carácter normativo de la gramática tradicional, el estructuralismo, a
partir del concepto central de \termino{sistema}, aportó un modelo consistente
para la descripción de las lenguas.
\end{destacado}

Dentro del estructuralismo, además de conceptos fundamentales para la
disciplina, como \textit{signo, sistema, lengua/habla, diacronía/sincronía,
significado/significante}, también se desarrollaron nociones clave, como las de
\termino{fonema} y \termino{morfema},\footnote{La noción de fonema se trata en
el capítulo~3 y la de morfema, en el capítulo~4.} que se incorporaron al
análisis junto a las unidades tradicionales, como la palabra y la oración.
Asimismo, para distinguir las «partes de la oración» o clases de palabras, se
buscó superar las distinciones basadas en conceptos semántico-nocionales y
establecer criterios más nítidos, asentados en pautas formales. Se puede
ejemplificar esta diferencia con las dos definiciones del verbo siguientes:

\begin{exe}
  \ex\label{ex:01-defverbo}
    \begin{xlist}
      \ex\label{ex:01-defverbo-a} palabra que manifiesta \sig{acción}.
      \ex\label{ex:01-defverbo-b} palabra que concuerda con el sustantivo núcleo del sujeto en
          número y persona.
    \end{xlist}
\end{exe}

La definición \exref{ex:01-defverbo-a}, que descansa en un criterio
semántico-nocional como el concepto de \sig{acción}, presenta la dificultad de
que no es aplicable a todos los verbos, ya que además de los que efectivamente
manifiestan una \sig{acción}\footnote{\label{fn:01-accion}Los verbos de
\sig{acción} responden a la pregunta \textit{¿Qué hizo X?} Así, frente a
\textit{¿Qué hizo Juan?} podemos responder: \textit{Juan corrió, comió} o
\textit{estudió}, pero no \textit{Juan creció} o \textit{nació}, que son verbos
de proceso, ni tampoco \textit{Juan tuvo fiebre} o \textit{era alto}, que son
verbos de estado (v. capítulo~6).} ---\forma{correr, comer, estudiar}---,
existen otros que señalan \sig{procesos} ---\forma{crecer, nacer} o
\forma{florecer}---, y otros que refieren \sig{estados} ---\forma{tener, ser,
durar}---. En cambio, \exref{ex:01-defverbo-b}, basada en un criterio formal,
como la concordancia en número y persona entre el verbo y su sujeto, resulta
aplicable en todos los casos.

En relación con la sintaxis, que fue adquiriendo cada vez más espacio frente a
la posición a la que la relegaba la gramática tradicional, el concepto de
\termino{constituyente inmediato} (CI) (v. capítulo~5), propuesto por el
estructuralista norteamericano Rulon \textcite{wells1947}, se convirtió en la
base del análisis sintáctico (ingl. \textit{parsing}), tal como lo conocemos
hoy en día.

Aunque el estructuralismo aportó un modelo completo para el estudio de la
lengua y un fundamento para el análisis sintáctico, con el correr del tiempo su
enfoque se consideró reduccionista porque, a partir de la distinción
saussureana entre \termino{lengua} y \termino{habla}, se dejó de lado el uso y
se puso el énfasis en el sistema. Este enfoque también recibió críticas por
establecer un modelo estático, que se centraba en el aspecto sincrónico y
dejaba de lado el cambio y la variación, aspectos esenciales de la dinámica de
las lenguas. Finalmente, también se lo cuestionó por ser fundamentalmente
taxonómico, es decir, enfocarse en la clasificación de las unidades para su
reconocimiento y etiquetado, dejando de lado su funcionamiento real en los
textos.

\subsection{La llegada del enfoque comunicativo}

A fin de subsanar los problemas que acabamos de mencionar, a mediados de los
años 80 del siglo pasado, se introdujo el denominado enfoque comunicativo,
surgido en la enseñanza de lenguas segundas y extranjeras. El objetivo de su
incorporación en el estudio de la lengua materna era subsanar la simplificación
y el alejamiento de los textos que el modelo estructural ofrecía. Se
incorporaron, entonces, propuestas provenientes de la gramática textual, del
análisis del discurso y de las teorías de la enunciación y de los actos de
habla, a las que se completaba con planteos tomados de los enfoques
predominantes en la enseñanza de lenguas extranjeras.

Asimismo, buscando remediar las limitaciones de la gramática oracional, se
introdujeron conceptos de la propuesta funcional\footnote{\label{fn:01-funcionalismo}El
funcionalismo parte de lo afirmado por los lingüistas estructuralistas de la
escuela de Praga (1929): «la lengua es un sistema de medios de expresión
apropiados para un fin» (\textit{TCLP}, I; \cite{aavv1929}). Los
funcionalistas conciben el lenguaje principalmente como un sistema de
comunicación, que se crea en la interacción social.} de
\textcite{halliday1976}. Se incorporaron en las aulas y manuales nociones
referidas al análisis textual, como coherencia, cohesión, progresión temática,
registro, entre otras. Sin embargo, como plantea \textcite[104]{gaiser2011},
«lo que el discurso de la didáctica de la lengua olvidó, es que, como dice el
propio Halliday, para aplicar a un texto la gramática funcional, se debe tener
una sólida base en gramática oracional».

En línea con lo señalado para Argentina, \textcite[18]{munguia2015} hace la
siguiente evaluación respecto de la aplicación del enfoque comunicativo en
México:

\begin{cita}
esta nueva reforma trató de enseñar «haciendo», es decir, los alumnos
aprenderían a escribir «escribiendo», a leer «leyendo» textos cuyas temáticas
fueran cercanas a sus intereses y motivaciones. Entonces, se prescindió de toda
teorización, por ello se eliminaron casi todos los temas relacionados con
gramática [\ldots]

Pero ya transcurrieron más de 20 años y los resultados siguen siendo muy
desalentadores [\ldots]. Las diferentes evaluaciones aplicadas muestran que el
enfoque comunicativo está lejos de haber logrado sus fines.
\end{cita}

En el ámbito local, lo que se observa en la actualidad es que, aunque la
propuesta comunicativa no cuajó en un modelo eficaz para la enseñanza de la
lengua materna, lo que logró fue arrinconar los contenidos gramaticales hasta
llevarlos a su casi completa desaparición de las aulas. En cuanto a la
sintaxis, o bien se ignora su tratamiento o bien se aplica el modelo más
conocido y que reproduce la mayoría de los manuales, que es el de cuño
estructural. Asimismo, la práctica tiende a ser realizada de un modo
predominantemente clasificatorio, sin establecer una genuina reflexión entre
los constituyentes que se reconocen en el análisis y los valores semánticos y
comunicativos que transmiten. Así, como expresa \textcite{crystal2007},

\begin{cita}
[n]o es sorprendente que la mayoría de la gente que aprendió análisis
sintáctico en el colegio termine sin poder ver la utilidad de este ejercicio, y
solo pueda recordar la gramática como una disciplina muerta e irrelevante, la
anatomía forense del lenguaje.
\end{cita}

\noindent Sin embargo, el mismo \textcite{crystal2007} añade que

\begin{cita}
[l]a realidad dice lo contrario. Las técnicas del análisis gramatical pueden
ser usadas para demostrar el enorme poder creativo del lenguaje ---cómo a
partir de un conjunto finito de patrones gramaticales, aún un niño pequeño
puede producir un conjunto infinito de oraciones--- (traducción nuestra).
\end{cita}

\subsection{Los aportes de la lingüística generativa}\label{sec:01-generativa}

La propuesta de la corriente conocida como gramática
generativa,\footnote{La denominación de la teoría se vincula con la concepción
de gramática que propone \textcite{chomsky1974a}: «la gramática de una lengua
será, por consiguiente, un dispositivo que genera todas las secuencias
gramaticales de la lengua y ninguna de las agramaticales» (v. capítulo~2).}
basada en los planteos de \textcite{chomsky1970,chomsky1974a}, en un sentido
representa una continuidad con el estructuralismo en lo referido a la
preocupación por el rigor metodológico, que en el caso de Chomsky lo lleva a
buscar una descripción del lenguaje cercana a las formulaciones lógicas y
matemáticas. Pero, en otro sentido, constituye una fuerte ruptura con todas las
teorías anteriores. A partir de los trabajos de Chomsky desde los años 60 del
siglo pasado, se produjo una revolución en el campo de la lingüística y las
ciencias cognitivas:\footnote{La lingüística integra el espacio de las ciencias
cognitivas junto a otras disciplinas interesadas en la exploración de la mente
(psicología, antropología, filosofía, inteligencia artificial, neurología,
etc.).}

\begin{cita}
Frente a los posicionamientos teóricos propios del estructuralismo europeo y
americano que concebían el lenguaje como institución social [\ldots], surge una
visión del lenguaje como facultad, que forma parte de la
mente/cerebro\footnote{Esta dupla hace referencia a la parte psíquica (mente) y
a su soporte físico (cerebro).} y debe ser aprehendido con técnicas y métodos
similares a los empleados en las ciencias naturales\ldots{}
\parencite[54]{mairal2012}.
\end{cita}

El lenguaje deja de considerarse un producto exclusivamente cultural y pasa a
estudiarse como una \termino{facultad}, en sentido biológico, una capacidad
propia y exclusiva de nuestra especie, con la que estamos dotados
naturalmente. Así, para \textcite[14]{chomsky1992},

\begin{cita}
[e]n el habla normal, uno no repite mecánicamente lo que ha oído, sino que
produce formas lingüísticas nuevas ---a menudo nuevas en la experiencia de uno
o incluso en la historia de la lengua--- y no hay límites para dicha
innovación.
\end{cita}

Pero ¿cómo es que esto se produce? La respuesta se vincula con lo que Chomsky
denominó \termino{el problema de Descartes} o \termino{de la infinitud
discreta}, que se refiere a la productividad del lenguaje, ya que, como
explican \textcite[17]{eguren2004}: «Los hablantes son capaces de producir un
número ilimitado de oraciones nuevas porque conocen un número finito de
unidades lingüísticas y saben, además, cómo combinarlas». Esta perspectiva
lleva a Chomsky a proponer que el lenguaje humano es recursivo\footnote{Un
proceso recursivo es aquel que se puede volver a aplicar al resultado de su
aplicación previa, en forma potencialmente indefinida. En gramática, la
subordinación es un proceso recursivo porque dentro de una oración se pueden
incrustar otras. Así en una oración simple como [\textsubscript{O} Tengo una
casa], podemos incluir otra, como en [\textsubscript{O1} Tengo una casa
[\textsubscript{O2} que está alejada del centro de la ciudad]], y conformar
una oración compuesta (v. capítulo~13).} y a plantear que la recursividad es
su característica más prominente, ya que la combinación de elementos finitos
puede dar lugar a un número infinito de oraciones.

El otro dilema que Chomsky ha planteado es el que se conoce como \termino{el
problema de Platón} o \termino{de la pobreza de los estímulos}, y tiene que ver
con otra importante cuestión abordada por esta teoría: ¿cómo se adquiere el
lenguaje? A diferencia de lo que sostenían los psicólogos
conductistas,\footnote{La psicología conductista, que tuvo gran influencia en
el estructuralismo norteamericano, considera que las conductas se producen
como respuesta a los estímulos del ambiente.} para quienes el lenguaje es una
conducta aprendida mediante ensayo y error, Chomsky sostiene que el lenguaje se
adquiere. ¿Pero en qué consiste la diferencia? Dentro de la mente hay
actividades cognitivas diferentes, con características y funciones
particulares. \textcite{chomsky1993} distingue entre «tareas naturales», que se
adquieren sin esfuerzo ---respirar, caminar, hablar--- y otras más complejas,
que tienen que ser aprendidas con empeño ---tocar el piano, jugar al ajedrez,
escribir, o construir teorías científicas---.

Sin embargo, aunque es natural, el desarrollo del lenguaje requiere del
estímulo que proviene de escucharlo en el medio ambiente, de la misma manera
que una semilla necesita la tierra y el agua para germinar en una planta:

\begin{cita}
[C]ualquier organismo necesita un entorno rico y estimulante para que sus
capacidades naturales emerjan. Volviendo a la imagen de que enseñar era como
permitir crecer bien a una flor, si no se riega la planta no va a crecer hasta
dar flor. No es que aprenda del agua a ser flor ---si fuera un árbol usaría la
misma agua para crecer hasta convertirse en árbol---. Pienso que sucede más o
menos lo mismo en el desarrollo humano, sin excluir el desarrollo del lenguaje
y del pensamiento \parencite[160]{chomsky1992}.
\end{cita}

Pero ¿qué implica esta nueva perspectiva que propone la gramática generativa
para el estudio del lenguaje? Si el lenguaje es una facultad alojada en
nuestra mente, se puede decir que cada hablante nativo tiene internalizada la
gramática de su lengua, en tanto puede identificar, formar e interpretar las
palabras y oraciones que se producen en ella. A este conocimiento tácito o no
consciente que el hablante posee acerca de su lengua, \textcite{chomsky1970}
lo denomina \termino{competencia} y lo diferencia de la \termino{actuación} o
uso real de la lengua en situaciones concretas.

Dado que el lenguaje es un conocimiento, \textcite{chomsky1989} plantea un
programa de investigación en torno a tres preguntas básicas:

\begin{itemize}
  \item ¿Qué es lo que constituye el conocimiento del lenguaje?
  \item ¿Cómo se adquiere ese conocimiento?
  \item ¿Cómo se utiliza?
\end{itemize}

\begin{ejercicio}\label{ej:01-facultad}
Teniendo en cuenta lo que se ha planteado en este capítulo y después de leer
muy atentamente el texto siguiente, complete lo que se pide más abajo, según
se muestra en el ejemplo resuelto~1.

\begin{exe}
  \ex El lenguaje es una facultad porque \respuesta{es una capacidad propia de
      nuestra especie}.\hfill{\footnotesize\sffamily\color{acento2}ejemplo resuelto}
  \ex Se trata de una facultad universal porque\completar
  \ex Es privativa de nuestra especie porque\completar
  \ex El lenguaje sirve para la comunicación, pero también cumple otros fines,
      como ser\completar
  \ex La forma oral del lenguaje es natural para el hombre porque\completar
  \ex Desde esta perspectiva, la escritura puede considerarse\completar
  \ex Sin el lenguaje la vida social no sería posible porque\completar
  \ex El lenguaje es parte de la dotación biológica de nuestra especie
      porque\completar
  \ex Para el desarrollo del lenguaje es necesario interactuar con el medio
      porque\completar
  \ex Somos capaces de producir combinaciones de unidades que constituyen
      nuevas oraciones que jamás hemos oído porque\completar
\end{exe}
\end{ejercicio}

\subsubsection{De la facultad del lenguaje a la gramática}

Para Chomsky la facultad del lenguaje consiste en \termino{principios} muy
generales que están presentes en todas las lenguas y constituyen la
\termino{gramática universal} (GU), de naturaleza mental o cognitiva. Estos
principios son comunes a todos los individuos de la especie y hacen posible que
los hablantes adquieran su lengua, es decir, que desarrollen la competencia
gramatical en el idioma hablado en su entorno.

Ahora bien, si los principios están en todas las lenguas, ¿por qué estas
difieren tanto? La respuesta de la teoría es que en cada lengua los principios
se realizan mediante \termino{parámetros}, que son los que establecen las
diferencias que se reconocen entre las lenguas.

Para explicar las diferencias, Chomsky recurre a la comparación con un
conmutador, cuyas complejas redes representan los distintos principios de la
facultad del lenguaje. Para que el sistema funcione, la red posee distintos
interruptores que deben fijarse en una de las dos posiciones disponibles para
cada uno. Cada posición adoptada por un interruptor representa el parámetro que
la lengua en cuestión adopta para cada principio y el tener que optar entre dos
posibilidades (y no entre infinitas) facilita a los niños la adquisición de la
propia lengua.

\figurapendiente{Una metáfora técnica}{fig:01-metafora}

Ahora veamos un sencillo ejemplo que permitirá ilustrar mejor la diferencia
entre principios y parámetros. Si observamos las siguientes oraciones, vemos
que mientras el español \exref{ex:01-sujnulo-a} y el italiano
\exref{ex:01-sujnulo-b} permiten omitir el sujeto,\footnote{Estas nociones
sobre la omisión del sujeto se amplían en el capítulo~6.} en otras lenguas,
como el inglés \exref{ex:01-sujnulo-c} y el francés \exref{ex:01-sujnulo-d},
una oración sin sujeto explícito está mal formada, es decir que resulta
agramatical.\footnote{El asterisco señala que se trata de una oración mal
formada o agramatical. El concepto de agramaticalidad se desarrolla en el
capítulo~2.}

\begin{exe}
  \ex\label{ex:01-sujnulo}
    \begin{xlist}
      \ex\label{ex:01-sujnulo-a} Comen perdices.
      \ex\label{ex:01-sujnulo-b} Mangiano le pernici.
      \ex\label{ex:01-sujnulo-c} \juicio{*}{Eat partridges} / They eat partridges.
      \ex\label{ex:01-sujnulo-d} \juicio{*}{Mangent des perdrix} / Ils mangent des perdrix.
    \end{xlist}
\end{exe}

¿Cómo podemos explicar estas diferencias? La variación que se observa en
\exref{ex:01-sujnulo} radica en la diferente forma en que se adopta el
principio que establece que todas las oraciones poseen un sujeto del cual se
predica. Como todo principio, este también tiene dos parámetros posibles de
realización (dos posiciones para el interruptor, siguiendo la metáfora
introducida arriba).

\begin{itemize}
  \item \textit{Opción 1}. El sujeto está presente obligatoriamente en todas
        las oraciones como sucede, por ejemplo, en inglés y en francés.
  \item \textit{Opción 2}. El sujeto puede omitirse, como, por ejemplo, en
        español y en italiano. Esta opción se conoce como \termino{parámetro
        del sujeto nulo} y las lenguas que permiten la omisión se denominan
        \textit{pro-drop} (con caída del pronombre sujeto).
\end{itemize}

\begin{recuadro}[Para recordar]
Los \termino{principios} son propiedades que tienen aplicación en todas las
lenguas, es decir que carecen de excepciones y constituyen la facultad misma
del lenguaje. Los \termino{parámetros} son las formas de realización de los
principios en las diferentes lenguas. La gramática de cada lengua no es más que
el conjunto de valores paramétricos que esa lengua adopta.
\end{recuadro}

\section{La estructura de la gramática}

Desde esta perspectiva la \termino{gramática} de una lengua es el mecanismo
cognitivo específico que permite formar adecuadamente las expresiones
lingüísticas encargadas de transmitir significados a través de sonidos. Y, dado
que son las estructuras del \termino{componente sintáctico}, al que
\textcite{chomsky1999} denomina el \termino{sistema computacional central}, las
encargadas de establecer la articulación entre ambos dominios ---el físico del
sonido y el conceptual del significado---, se explica la centralidad otorgada a
este componente en el esquema de gramática que plantea. La gramática también
posee \termino{dos componentes de interfaz}, \termino{forma fonética} y
\termino{forma lógica}, encargados, respectivamente, de determinar la
interpretación fonológica y semántica de las cadenas generadas por la sintaxis,
para lo cual conectan las representaciones abstractas producidas, con el
sistema articulatorio-perceptivo (que emite y percibe los sonidos) y con el
sistema conceptual-intencional (que interpreta y comunica los
significados).\footnote{El modelo presentado en \textcite{hauser2002}
distingue entre: a) la Facultad del Lenguaje en sentido amplio (FLA), que
incluye los sistemas «periféricos» de procesamiento (articulatorio-perceptivo)
y de conceptualización (conceptual-intencional); y b) la Facultad del Lenguaje
en sentido estricto (FLE), cuyo componente central es el dispositivo
computacional, dotado de la que Chomsky ha señalado como la principal
innovación del lenguaje: la recursividad.} El esquema de la gramática se
completa con el componente léxico ---también conocido como \termino{diccionario}
o \termino{lexicón mental}---, que almacena y proporciona las unidades léxicas
que necesita el componente sintáctico para alimentar las estructuras
lingüísticas que produce. La figura~\ref{fig:01-modelo} muestra el esquema
global de la gramática, según el modelo presentado en \textcite{chomsky1999}:

\begin{figure}[htbp]
  \centering
  \begin{tikzpicture}[
      font=\sffamily\small,
      caja/.style={draw=acento, thick, fill=acentoclaro, align=center,
                   minimum width=34mm, minimum height=9mm, inner sep=3pt},
      externo/.style={draw=gristexto, dashed, align=center,
                      minimum width=34mm, minimum height=9mm, inner sep=3pt,
                      text=gristexto},
      flecha/.style={-{Stealth[length=2.2mm]}, thick, acento},
      enlace/.style={{Stealth[length=2mm]}-{Stealth[length=2mm]}, gristexto}]
    \node[caja] (lex) at (0,0) {Léxico\\{\footnotesize (diccionario o lexicón mental)}};
    \node[caja] (sint) at (0,-1.9) {Sistema computacional\\{\footnotesize (componente sintáctico)}};
    \node[caja] (ff) at (-2.6,-3.9) {Forma fonética};
    \node[caja] (fl) at (2.6,-3.9) {Forma lógica};
    \node[externo] (ap) at (-2.6,-5.6) {Sistema\\articulatorio-perceptivo};
    \node[externo] (ci) at (2.6,-5.6) {Sistema\\conceptual-intencional};
    \draw[flecha] (lex) -- (sint);
    \draw[flecha] (sint) -- (ff);
    \draw[flecha] (sint) -- (fl);
    \draw[enlace] (ff) -- (ap);
    \draw[enlace] (fl) -- (ci);
  \end{tikzpicture}
  \caption{El modelo de la gramática}\label{fig:01-modelo}
\end{figure}

Si bien desde la perspectiva de la propuesta de Chomsky resulta perfectamente
entendible la centralidad otorgada a la sintaxis, en numerosas ocasiones esta
postura ha sido criticada por lingüistas que ponen el acento en los aspectos
semánticos\footnote{La corriente conocida como Semántica Generativa, que se
desarrolló entre 1968 y 1975, con representantes como G.~Lakoff, J.~Ross,
J.~McCawley y Ch.~Fillmore, entre otros lingüistas formados dentro del tronco
generativo, cuestionaba la centralidad de la sintaxis.} o en el valor
comunicativo de las combinaciones. Los cuestionamientos recibidos han ido
llevando a la teoría a ir otorgando cada vez más relevancia a las cuestiones
léxico-semánticas.\footnote{En \textit{Observaciones sobre la nominalización}
\parencite{chomsky1974b}, como respuesta a los cuestionamientos de los
semánticos generativos, Chomsky propuso la hipótesis lexicalista, que otorga
mayor prominencia a la palabra.} En este sentido ha sido fundamental el
\termino{principio de proyección} propuesto por \textcite{chomsky1981}, el cual
plantea que las estructuras sintácticas se proyectan a partir de las unidades
léxicas. Pero ¿qué implicaciones tiene este principio en la combinatoria de las
palabras?

\subsection{La proyección sintagmática de las palabras}\label{sec:01-proyeccion}

Para \textcite[21]{bosque2009} «podemos al abordar el análisis sintáctico
\enquote*{empezar por las oraciones}, como se hacía en la tradición, o bien
\enquote*{empezar por las palabras}, según se recomienda ahora». Los
mencionados autores se pronuncian por esta segunda opción fundados en que las
que se van combinando entre sí son las palabras, primero «en grupos pequeños y
luego en secuencias cada vez más complejas». Esos «grupos» en que se combinan
las palabras son las denominadas unidades intermedias de la gramática, o sea
los \termino{sintagmas}, que son la máxima proyección o expansión de un núcleo.
Y es el núcleo, según la clase de palabra a la que pertenece ---nombre (N),
verbo (V), adjetivo (A), adverbio (Adv), preposición (P)\footnote{\label{fn:01-lexicas}Nos
limitamos aquí a la proyección de las clases de palabras léxicas, que son las
que poseen contenido descriptivo inherente: \textit{árbol} (N), \textit{viejo}
(A), \textit{rápidamente} (Adv), \textit{correr} (V), \textit{para} (P). Las
palabras léxicas se diferencian de las palabras funcionales, como los
determinativos (\textit{el, este, un, mi}, etc.) y las conjunciones
(\textit{y, pero, o, que}, etc.), cuyo significado se relaciona con la función
gramatical que cumplen (v. \cite{giammatteo2009}).}---, el que le da su clase o
categoría al sintagma (SN, SV, SA, SAdv, SP). A su vez, los grupos así
conformados se seguirán recombinando para ir conformando secuencias cada vez
mayores, como \exref{ex:01-roma-a}, representada también en el diagrama
arbóreo \exref{ex:01-roma-b}:

\begin{exe}
  \ex\label{ex:01-roma}
    \begin{xlist}
      \ex\label{ex:01-roma-a} Francisco vivía en Roma.
      \ex\label{ex:01-roma-b}
        \begin{forest} gramatica
          [O
            [SN [Francisco, roof]]
            [SV
              [V [vivía]]
              [SP [P [en]] [SN [Roma, roof]]]]]
        \end{forest}
    \end{xlist}
\end{exe}

Dado que en la oración los elementos se van combinando desde el interior, es
decir, de derecha a izquierda, la lectura de \exref{ex:01-roma-b} procede de
abajo hacia arriba: en primer lugar, la P \forma{en} se une con su complemento
\forma{Roma}, formando un SP ---\forma{en Roma}---; luego ese SP se combina con
la forma verbal \forma{vive} y conforma el SV \forma{vive en Roma}, que, a su
vez, se combina con el SN \forma{Francisco} formando una oración. Asimismo,
dado que \forma{vive} concuerda con el SN \forma{Francisco} en persona y número
(ambos están en 3.ª persona singular), este último hará las veces de sujeto de
la oración resultante.\pendiente{El ejemplo dice \textit{vivía}; el texto,
\textit{vive}}

Ahora bien, ¿cualquier combinación de elementos léxicos es válida? ¿Qué pasa si
combinamos el SP \forma{en Roma} con el verbo \forma{nieva} y, al SV resultante
\forma{nieva en Roma}, lo unimos con el SN \forma{Francisco}? Si bien los SS
combinados siguen siendo de la misma clase de palabras, la oración
\juicio{*}{\forma{Francisco nieva en Roma}} no estaría bien formada porque
\forma{nevar} es un V impersonal que no admite un sujeto léxico como
\forma{Francisco}.

\begin{ejercicio}\label{ej:01-combinaciones}
Considere las siguientes combinaciones e intente explicar por qué resultan
anómalas o bajo qué condiciones podrían ser aceptables, según se muestra en el
ejemplo resuelto~1.

\begin{exe}
  \ex Francisco maúlla en Roma.
  \ex La lluvia vive en Roma.
  \ex Francisco tiene en Roma.
  \ex Francisco vive en triste.
  \ex Francisco reza la piedra.
  \ex La piedra reza en Francisco.
\end{exe}

\begin{resuelto}
1. El V \forma{maullar} exige un sujeto con los rasgos semánticos
\sig{animado} y \sig{felino}. La O sería aceptable si \forma{Francisco}
remitiera a un individuo con esas características.
\end{resuelto}
\end{ejercicio}

Los casos del ejercicio~\ref{ej:01-combinaciones} muestran que las unidades
léxicas tienen exigencias combinatorias tanto sintácticas como semánticas. En
cuanto a los requisitos sintácticos, vemos que la P requiere un SN como
complemento y que un V transitivo como \forma{tener} exige un complemento
objeto directo (COD) (v. capítulo~7). Respecto de las exigencias semánticas:
\forma{maullar} pide un sujeto animado \forma{gato} (o algún otro felino);
\forma{rezar} exige un COD del tipo \forma{oración} o \forma{plegaria}; la P
\forma{en} requiere no solo un SN como complemento, sino que este debe referir
a un lugar ---\forma{en Roma, en la colina}--- y también, con un V como
\forma{rezar}, puede indicar una lengua determinada ---\forma{Francisco reza en
latín}---. Cuestiones de este tipo, entre otras muchas en las que no podemos
extendernos aquí, pero que se irán desarrollando a lo largo de este libro,
hacen que la gramática no pueda quedarse exclusivamente en el estudio de las
formas, sino que también deba atender a los significados.

\begin{destacado}
En palabras de \textcite[13]{bosque2009}, «el análisis de la combinatoria
gramatical no puede \enquote*{poner entre paréntesis} el significado, sino más
bien todo lo contrario: es tarea esencial de la sintaxis explicar lo que
significan las combinaciones de palabras».
\end{destacado}

\subsection{Las palabras en contexto}

Ahora bien, ¿es lo que acabamos de plantear suficiente para dar cuenta
acabadamente de lo que manifiestan las expresiones lingüísticas? Veamos un
nuevo ejemplo en contexto:

\begin{exe}
  \ex\label{ex:01-francisco} Francisco vivía en Roma, en la casa de Santa
      Marta, un edificio adyacente a la Basílica de San Pedro. Durante un
      tiempo, en la misma Ciudad del Vaticano, convivió con su predecesor, el
      papa emérito, Benedicto XVI, que había fallecido en 2022.
\end{exe}

Aunque el sujeto y el verbo de la primera oración de \exref{ex:01-francisco}
siguen siendo los mismos de \exref{ex:01-roma}, nuestro conocimiento de mundo
---la enciclopedia que cada hablante tiene incorporada en su mente según sus
experiencias de vida--- y las pistas contextuales que aportan los elementos
léxicos ---\forma{la casa de Santa Marta, la Basílica de San Pedro, la Ciudad
del Vaticano, el papa emérito Benedicto XVI}--- nos llevan a asociar el nombre
\forma{Francisco} con el del papa fallecido en 2025. Si no fuéramos capaces de
establecer esta vinculación, nuestra interpretación de \exref{ex:01-francisco}
quedaría incompleta. Así, si bien dijimos que la perspectiva comunicativa que
excluye a la gramática no resultó eficaz para la enseñanza, esto no significa
que debamos desentendernos completamente del funcionamiento de las estructuras
en contexto. Como sostiene \textcite[s/p]{brucart2018},

\begin{cita}
[n]o integrar la gramática oracional en una perspectiva comunicativa equivale a
diseñar una ciudad con la trama urbana y los servicios preparados, pero sin
edificios ni principios constructivos.
\end{cita}

\subsection{Gramática y comprensión}

Si queremos integrar la gramática en la comprensión textual, es necesario
revisar las distintas concepciones acerca de la lectura a lo largo de la
historia. Según plantea \textcite{pozo2014}, en 5000 años de desarrollo, la
lectura ha ido cambiando, por lo que reconoce tres etapas:

\begin{itemize}
  \item La de \termino{lectura reproductiva o literal}, propia de los primeros
        tiempos, cuando aprender a leer era repetir y la función del lector no
        era encontrar el sentido de lo que leía, sino retener el texto
        almacenado en su mente.
  \item La de la \termino{lectura escolástica}, hacia el fin de la Edad Media,
        cuando la función de la lectura ya no era repetir lo que decía el
        texto, sino reproducir una interpretación autorizada de él.
  \item La de la \termino{lectura moderna}, que surge a partir del
        Renacimiento con la invención de la imprenta, cuando aparece el lector
        crítico o analítico. Se trata de un lector que va a interpretar,
        analizar, criticar, dar sentido y relacionarse epistémicamente con el
        texto.
\end{itemize}

A las tres etapas señaladas por Pozo, hoy día ya es necesario incorporar una
cuarta: la de \termino{la lectura hipertextual}, caracterizada por la
multimodalidad (que conecta con otros medios), la interactividad (que vincula
con otros usuarios), la flexibilidad (que permite otros recorridos) y la
intertextualidad (que relaciona con otros textos).

\figurapendiente{La lectura hipertextual}{fig:01-hipertextual}

En la actualidad, la comprensión lectora se considera un proceso formado por
múltiples procedimientos cognitivos que posibilitan al lector modificar su
pensamiento respecto de la realidad mediante la información que descubre en el
texto. Por tanto, lo que interesa en la enseñanza es formar lectores críticos
capaces de interpelar al texto, de interactuar con él y ser capaces de
relacionarlo con otros textos, para lo cual la simple habilidad de decodificar
grafemas\footnote{Los grafemas son las unidades mínimas de la escritura de una
lengua (\textit{DLE}, en línea).} no alcanza. Se requiere una alfabetización
integral que sume otros saberes «lectores» ---los aspectos enunciativos, los
géneros discursivos, los factores histórico-culturales de producción y
recepción---, pero sin excluir un conocimiento gramatical explícito, que ayude
a reconocer las «pistas» que ofrecen los textos para su cabal interpretación.

Pero ¿cómo se puede vincular las estructuras ---motor de la lengua--- con los
significados, que les dan sentido y, a su vez estos, con los valores
comunicativos que los hablantes les adjudican en sus distintos usos? Volvamos a
recurrir a una historieta. La viñeta de abajo muestra tres personajes: el
hablante, un niño aborigen; su interlocutora, a la que se dirige con el
apelativo familiar \forma{ma}; y un tercer personaje que lleva un traje de la
época de la conquista, a quien podríamos identificar con Cristóbal Colón. Pero
¿qué dice el texto que acompaña la imagen?

\figurapendiente{Fuente:
\url{https://m.facebook.com/866981573634920/photos/a.937101296622947/1274222342910839/?type=3}}{fig:01-colon}

Aunque breve, podemos decir que el texto de la viñeta constituye una
argumentación completa. Sin pretender un análisis exhaustivo, señalaremos
algunos de los principales recursos gramaticales que emplea el texto para
lograr su efectividad.

\begin{itemize}
  \item En la primera oración, el niño se vale del discurso referido
        indirecto\footnote{A diferencia del discurso referido directo, en el
        que el hablante reproduce literalmente lo dicho por otro ---p.~ej.
        \textit{Juan dijo: Ahora no puedo ir}---, en el indirecto el hablante
        ajusta lo dicho a su situación comunicativa ---\textit{Juan dijo que en
        ese momento no podía venir}---.} como estrategia\footnote{Se trata de un
        recurso evidencial mediante el cual el hablante expresa la fuente de la
        que deriva la información. El tema de modalidad y evidencialidad se
        desarrolla en el capítulo~14.} que le permite reproducir lo dicho por el
        conquistador ---\forma{hay un señor que dice que nos descubrió}--- y, al
        mismo tiempo, tomar distancia y desligar su responsabilidad respecto de
        un argumento con el que no concuerda.
  \item En el inicio de la segunda oración, hay una palabra
        funcional,\footnote{V. nota~\ref{fn:01-lexicas}.} la conjunción
        adversativa \forma{pero}, que conecta con la oración anterior e
        introduce el contraargumento mediante el cual el niño refuta lo dicho
        por el conquistador dejando en evidencia que no estaban dadas las
        condiciones para «ser descubiertos»,\footnote{La acepción 3 de
        \textit{descubrir} en el \textit{DLE}, en línea, dice: «Hallar lo que
        estaba ignorado o escondido, principalmente tierras o mares
        desconocidos».} ya que «no estaban escondidos».
  \item En la tercera oración, el adverbio deíctico \forma{acá} seguido del SP
        \forma{en casa}, que lo desarrolla,\footnote{Constituye una aposición
        explicativa (v. capítulo~5).} remite al contexto de situación y le
        permite al niño reafirmarse no solo en el lugar de hablante, sino en el
        de «dueño de casa». En el contexto de la historieta, nuestro
        conocimiento de mundo nos lleva a inferir que tanto \forma{acá} como
        \forma{en casa} refieren a América, el territorio que el «señor» dice
        haber descubierto. Asimismo, \forma{en casa}, junto con el apelativo
        \forma{ma} de la primera oración, remiten a un contexto familiar que
        evidencia el carácter de intruso del conquistador y desbarata sus
        dichos y pretensiones.
  \item La oración final contiene la conclusión que, en línea con el trato
        respetuoso que el niño le dispensa al conquistador al referirse a él
        como \forma{señor}, también se manifiesta en un tono morigerado, aunque
        no exento de sarcasmo: \forma{me parece que él está perdido}. La ironía
        que sobrevuela todo el texto remata en una doble posibilidad
        interpretativa para \forma{estar perdido}:
        \begin{itemize}
          \item \forma{está perdido}\textsubscript{1} porque «ha errado el
                camino», ya que, como se sabe, Colón se dirigía a las Indias y
                no a América,
          \item \forma{está perdido}\textsubscript{2} porque «está confundido»
                y no es cierto lo que afirma.
        \end{itemize}
\end{itemize}

En suma, como venimos señalando, los marcos teóricos de corte comunicativo que
hacen hincapié en el uso, pero se desentienden del estudio formal de los
fenómenos no son suficientes para la enseñanza, pero tampoco lo son los
enfoques formales cuando no tienen en cuenta que la lengua es un dispositivo
mental que los hablantes empleamos con múltiples intenciones para innumerables
fines. En este sentido, la gramática ha de servir para que los conocimientos
sobre los recursos lingüísticos estudiados puedan proyectarse en los textos y
colaborar en su comprensión.

\section{La gramática como descubrimiento}

Según venimos viendo, la lengua constituye un mecanismo de gran complejidad
diseñado para transmitir, mediante el uso de sonidos distintivos (fonemas o
bien grafemas, si se trata de la escritura), información (es decir,
significados) sobre el mundo y sobre nosotros mismos, para lo cual se vale de
una estructura con diferentes componentes interconectados.

El planteo central de la propuesta que se desarrolla en este libro, y que aquí
estamos solo presentando, se asienta en una base formal que permite dar cuenta
de las restricciones y correlaciones que conforman la combinatoria lingüística.

\begin{destacado}
Desde esta perspectiva, estudiar gramática no es otra cosa que hacer explícito
ese saber operativo (\textit{know how}, saber hacer o conocimiento tácito) que
el hablante común posee y que le permite comprender y producir las infinitas
oraciones de su lengua.
\end{destacado}

Sin embargo, como sostienen \textcite[69]{bosque2016}, por lo general, no se
enseña «a observar, no tanto lo anómalo o lo inusitado, sino más bien lo
sorprendente cuando se halla dentro de lo cotidiano y es relevante para
comprender las generalizaciones». Por eso, en esta propuesta el énfasis está
puesto en llamar la atención sobre el lenguaje y sus particularidades, a fin de
comenzar a reflexionar sobre él.

El punto de partida será la propia competencia, lo que como hablante nativo
cada uno «sabe» acerca de su lengua.

El objetivo es hacer consciente (explícito) ese conocimiento que tenemos
internalizado para que pueda ser utilizado eficazmente en distintas
situaciones, y así

\begin{itemize}[label=$\rightarrow$]
  \item poder comprender mejor lo que nos dicen, y
  \item ser mejor comprendidos en nuestros intercambios comunicativos.
\end{itemize}

Este proceso que va del «saber hacer» al conocer consciente comprende dos
momentos:

\begin{itemize}
  \item \termino{el «extrañamiento»}, que implica una toma de distancia de la
        inmediatez que todos tenemos respecto de la propia lengua, para poder
        considerarla como un objeto de estudio, cuyo abordaje implica
        categorías y criterios de análisis específicos;
  \item \termino{el «descubrimiento»}, que lleva a reconocer no solo qué
        formas existen en la lengua y qué restricciones y correlaciones de
        formación presentan, sino también cuáles resultan posibles y cuáles
        imposibles y por qué. Asimismo, estas mismas constricciones permitirán
        entender los aparentes «desvíos» en que incurren los hablantes (v. los
        ejemplos de \exref{ex:01-mente-nombres}).
\end{itemize}

El enfoque que se propone, al que denominamos léxico-gramatical porque se
asienta en estos dos componentes esenciales de la estructura lingüística, tiene
en cuenta que son los elementos léxicos, con una determinada configuración
fonológica y morfosemántica, los que, al operar en relación sintagmática, van
conformando unidades significativas mayores cada vez más amplias y abarcativas
\parencite{giammatteo2017}. Teniendo en cuenta las unidades resultantes, hemos
organizado la propuesta en torno a cuatro ejes principales:

\begin{figure}[htbp]
  \centering
  \begin{tikzpicture}[
      font=\sffamily\small,
      eje/.style={draw=acento, thick, fill=acentoclaro, align=center,
                  text width=25mm, minimum height=10mm, inner sep=3pt},
      unid/.style={align=center, text width=25mm, font=\sffamily\footnotesize},
      flecha/.style={-{Stealth[length=2.2mm]}, thick, acento}]
    \node[eje] (e1) at (0,0)    {\textbf{1.\textsuperscript{er} eje}\\Nivel subléxico};
    \node[eje] (e2) at (3.2,0)  {\textbf{2.\textsuperscript{do} eje}\\Nivel léxico};
    \node[eje] (e3) at (6.4,0)  {\textbf{3.\textsuperscript{er} eje}\\Nivel sintagmático};
    \node[eje] (e4) at (9.6,0)  {\textbf{4.\textsuperscript{to} eje}\\Nivel textual};
    \node[unid, below=2mm of e1, anchor=north] {Fonemas\\Morfemas};
    \node[unid, below=2mm of e2, anchor=north] {Palabras};
    \node[unid, below=2mm of e3, anchor=north] {Sintagmas\\Oraciones};
    \node[unid, below=2mm of e4, anchor=north] {Párrafos\\Secuencias\\Textos};
    \draw[flecha] (e1) -- (e2);
    \draw[flecha] (e2) -- (e3);
    \draw[flecha] (e3) -- (e4);
  \end{tikzpicture}
  \caption{Ejes de la propuesta}\label{fig:01-ejes}
\end{figure}

En lo que sigue, se presentan, a modo de ejemplos, algunos casos que se podrían
considerar a fin de producir el «extrañamiento» respecto del uso cotidiano y
propiciar, a través de la reflexión, el «descubrimiento» de sus posibilidades e
imposibilidades combinatorias, lo cual debe conducir a reconocer el
funcionamiento de las estructuras y de sus significados, y a identificar sus
valores comunicativos y expresivos.

\subsection{Del interior de la palabra al reconocimiento de sus valores en contexto}

\subsubsection{Restricciones y correlaciones combinatorias en el nivel subléxico}\label{sec:01-sublexico}

En el nivel subléxico referido al sonido se encuentran los \termino{fonemas},
que son unidades fónicas sin significado en sí mismas, pero cuyo contraste
permite distinguir significados como, por ejemplo, entre
/ˈpaso/\footnote{En la transcripción fonológica, los fonemas convencionalmente
van entre barras inclinadas y el acento se coloca delante de la sílaba (v.
capítulo~3).} y /ˈpeso/ (la diferencia es la presencia de
/-a-/\footnote{Los guiones indican que una forma no es libre, sino que tiene
algún contexto antes /-a/ o después /a-/, o bien en ambas posiciones /-a-/.}
frente a /-e-/) o entre /ˈpeso/ y /ˈbeso/ (la diferencia se da entre /p-/ y
/b-/). En cuanto a la combinatoria de estas unidades, podemos decir que la
estructura fonológica del español admite algunas combinaciones y rechaza otras;
así, por ejemplo, es posible formar cadenas como \forma{\textbf{tr}aza} o
\forma{\textbf{tr}ama}, en las que interviene el grupo consonántico
\forma{tr-}, pero nunca podríamos encontrar \juicio{*}{\forma{\textbf{sr}aza}}
o \juicio{*}{\forma{\textbf{sr}ama}}, ya que la combinación \forma{sr-} no
constituye un patrón fonológico de nuestra lengua. A veces, para responder a los
esquemas propios de la lengua, se necesitan ajustes fonológicos: por ejemplo,
algunas palabras provenientes del latín o de otras lenguas que admiten grupos
consonánticos iniciales que en español no se aceptan, incorporan una /e-/
delante, para asimilarse a la norma de pronunciación del
español:\footnote{El signo > se lee: «se convierte en».}

\begin{exe}
  \ex\label{ex:01-eprotetica}
    \begin{tabular}[t]{@{}l l c l l@{}}
      \forma{schola}  & (lat.)  & > & \forma{e-scuela} (esp.) & (ingl. \forma{school}) \\
      \forma{stilus}  & (lat.)  & > & \forma{e-stilo} (esp.)  & (ingl. \forma{style}) \\
      \forma{stress}  & (ingl.) & > & \forma{e-strés} (esp.)  & \\
      \forma{scanner} & (ingl.) & > & \forma{e-scáner} (esp.) & \\
    \end{tabular}
\end{exe}

Al combinarse dentro de la palabra, los fonemas van formando unidades
significativas\footnote{Según se verá en el capítulo~4, el significado de los
morfemas puede ser léxico (de contenido) o puramente gramatical (funcional). En
ese capítulo también se desarrolla la noción de morfema y los conceptos
relacionados con la estructura interna de la palabra.} denominadas
\termino{morfemas}; así, a una cadena como /ˈutil/ (= \sig{que produce
beneficio}\footnote{Convencionalmente los significados se colocan entre comillas
simples.}) se le puede añadir un prefijo de significado negativo como /in-/
para formar /inˈutil/ (\sig{que no produce beneficio}) y luego, si añadimos el
sufijo \forma{-mente}, el resultado será /inˈutilmente/ (\sig{de manera inútil
o que no produce beneficio}). Pero también existen restricciones para la
combinación de morfemas: con el mismo mecanismo de \forma{inútilmente} también
podemos formar \forma{alegremente, rápidamente, infantilmente} (\sig{de manera
alegre, rápida o infantil}), pero no resultan posibles combinaciones como
\juicio{*}{\forma{vasomente}}, \juicio{*}{\forma{libromente}} o
\juicio{*}{\forma{mesamente}}, ya que el sufijo \forma{-mente} se puede unir a
adjetivos, como \forma{alegre, rápido} o \forma{infantil}, pero no a
sustantivos, como \forma{vaso, libro} o \forma{mesa}.

Estos conocimientos sobre la conformación de las palabras los tenemos
internalizados y, por lo general, no reflexionamos sobre ellos, los usamos
automáticamente y no dudamos en emplearlos cuando formamos nuevas palabras;
así, a partir del adjetivo coloquial \forma{trucho}, no titubeamos en formar el
derivado \forma{truchamente} (\sig{en forma trucha}, es decir, \sig{falsa o
ilegal}).

De este modo, según vamos viendo, el hilo conductor de la reflexión sobre la
lengua son los propios juicios acerca de la buena o mala formación de las
distintas construcciones, es decir, sobre su gramaticalidad o agramaticalidad.
Pero muchas veces estos conceptos también deberán ser complementados con la
noción de aceptabilidad que, en ciertos contextos, puede permitir extender los
límites de formación (v. capítulo~2), tal como muestran los ejemplos
siguientes, en los que, contrariamente a lo que hemos dicho más arriba, el
sufijo \forma{-mente} aparece unido a sustantivos:

\begin{exe}
  \ex\label{ex:01-mente-nombres}
    \begin{xlist}
      \ex Llegará, \forma{tanga\textbf{mente}}, mi muerte enamorada (H.~Ferrer,
          \textit{Balada para mi muerte}).
      \ex Zarpar en el barco sin mirar atrás, subirse a la patera que aguantará
          sólo mientras pueda,\ldots{} Las mayorías soportan
          \forma{golondrina\textbf{mente}} un presente atroz y aplastador. O
          terminan mansamente confinados como esclavos en talleres textiles
          \ldots{} (C.~Rafael, \textit{La tos de la
          esclavitud}).\pendiente{Falta el dominio del enlace en el original}
      \ex Muestra de esta insensibilidad es cuando por ejemplo estos bancos
          neoliberales provocan una y otra vez las olas especulativas en los
          tipos de cambio\ldots{} lo que les permite a través del Mercado de
          Divisas mover sus reservas y utilidades hacia donde
          \forma{golondrina\textbf{mente}} les conviene (A.~Tirado Pazos, 2016,
          \textit{La paradoja de México: Del neoliberalismo al liberalismo
          social},
          \url{https://www.ruizhealytimes.com/economia-y-negocios/la-paradoja-de-mexico-del-neoliberalismo-al-liberalismo-social}).
    \end{xlist}
\end{exe}

¿Qué explicación podemos ofrecer para estas «aparentes» anomalías? A diferencia
de \juicio{*}{\forma{vaso\textbf{mente}}}, \juicio{*}{\forma{libro\textbf{mente}}}
o \juicio{*}{\forma{mesa\textbf{mente}}}, en los que es imposible asignar
significado al conjunto (\juicio{*}{\sig{de manera vaso, libro o mesa}}), los
ejemplos de \exref{ex:01-mente-nombres}, en cambio, son aceptables para los
hablantes porque resultan interpretables en un contexto determinado. Así,
\forma{tangamente} significa \sig{de manera tanguera} y
\forma{golondrinamente} se entiende en el sentido de \sig{como se comportan las
golondrinas migratorias, en forma estacional}.

\begin{ejercicio}\label{ej:01-diccionariamente}
Lea atentamente el siguiente fragmento y responda si es o no aceptable la
palabra subrayada en el texto. Si la respuesta es afirmativa, explique cuál es
su significado.

\begin{cita}
Miller también fue pionero del estudio de cómo los chicos aprenden palabras,
no \subrayado{diccionariamente} sino en contexto: nadie puede estudiar desde el
diccionario (D.~Golombek, «Las palabras (y las cosas). Y los viajes por el
léxico», \textit{La Nación}, 22/01/2017).
\end{cita}
\end{ejercicio}

\subsubsection{Comparación y contraste: los pares mínimos}

Según estamos viendo, entonces, se parte de una focalización que lleva a
descubrir a través de ejemplos, las estructuras que subyacen a la lengua, lo
que luego podrá hacerse extensivo a otros casos similares para llegar, si cabe,
a establecer generalizaciones. El proceso cognitivo en que se basa «el
descubrimiento» es la comparación, que permite reconocer, a través del cotejo
de los que se conocen como \termino{pares mínimos}
\parencite{bosque2016,brucart2018}, qué comparten y en qué se diferencian las
construcciones examinadas.\footnote{En \textcite{bosque2016} se plantea que los
pares mínimos pueden ser dos secuencias sencillas y gramaticales que se
diferencian en un solo factor o dos secuencias, una de las cuales es
agramatical.} Aunque son muy utilizados en fonología, ya que sirven para
distinguir sonidos contrastivos como los fonemas
(v.~§\ref{sec:01-sublexico}), los pares mínimos no suelen emplearse para
establecer comparaciones en otros niveles de la estructura lingüística. Para
mostrar su utilidad, a modo de ejemplo, podemos comparar entre sí los
siguientes pares de palabras cuyas estructuras ---un prefijo seguido de una
base\footnote{El concepto de base es relacional: es de donde se parte para
hacer una derivación; por ejemplo, en \textit{inútil} la base es \textit{útil},
pero en \textit{inútilmente}, la base es \textit{inútil} (v. capítulo~4:
§6.1).}--- pueden parecer idénticas, aunque, como veremos, no lo son:

\begin{exe}
  \ex\label{ex:01-despistar} \forma{descentrar} / \forma{despistar}
  \ex\label{ex:01-autopista} \forma{autorretrato} / \forma{autopista}
\end{exe}

En \exref{ex:01-despistar}, ambas palabras están formadas por un sustantivo
respectivamente, \forma{centro} y \forma{pista}, a los que se les adjunta el
mismo prefijo \forma{des-}. Sin embargo, solo en la primera es posible derivar
un verbo del sustantivo \forma{centro} > \forma{centrar}, al que se le puede
añadir el prefijo \forma{des-}, que implica la realización de la \sig{acción
opuesta} \forma{centrar} ≠ \forma{descentrar}, y formar así una palabra
derivada. En \forma{despistar}, en cambio, no podemos partir de un inexistente
verbo \juicio{*}{\forma{pistar}}, lo que hace que \forma{despistar} provenga
directamente del nombre \forma{pista}, al que se le adjuntan, simultáneamente,
un prefijo y un sufijo, según el proceso de formación de palabras que se conoce
como parasíntesis (v. capítulo~4: §7.3).

\begin{exe}
  \ex\label{ex:01-parasintesis}
    \begin{tabular}[t]{@{}l@{\quad}l@{}}
      \forma{centro} > \forma{centr\textbf{ar}} > \forma{\textbf{des}centrar} & (palabra derivada) \\
      \forma{pista} > \forma{\textbf{des}pist\textbf{ar}}                     & (palabra parasintética) \\
    \end{tabular}
\end{exe}

En cuanto a los significados, ambas palabras constituyen el tipo de opuestos
que \textcite[18]{cruse2001} denomina \termino{reversos}: «pares tales como:
\textit{entrar: salir, subir: bajar, avanzar: retroceder, ascender: descender,
aparecer: desaparecer, unir: desunir, vestir: desvestir}. Todos ellos denotan
cambios de estado de algún tipo que son, en principio, reversibles». Desde esta
perspectiva, \forma{descentrar} es el proceso inverso de \forma{centrar}, o sea
«sacar del centro» y, de modo semejante, aunque \juicio{*}{\forma{pistar}} no
exista, \forma{despistar}, en una primera acepción se entiende como su opuesto:
«hacer perder la pista» (\textit{DLE}, en línea), pero, además,
\forma{despistar} adopta un valor derivado metafóricamente: «distraer,
desorientar o confundir» (\textit{DLE}, en línea), ya que quien se sale de
pista queda, por lo general, desorientado o confundido.

En cuanto a las palabras de \exref{ex:01-auto-estructura},
\forma{autorretrato}, formada por la unión del prefijo \forma{auto-} y el
sustantivo \forma{retrato}, se corresponde con el significado esperado de
\sig{retrato de uno mismo}. Pero si aplicamos el mismo procedimiento en
\forma{autopista}, nos lleva a un significado inexistente \juicio{\#}{\sig{pista
de uno mismo}}. En este segundo caso, como en otros muchos similares en español
(\forma{bicisenda, homofóbico},\footnote{Se trata también de palabras
compuestas con un acortamiento en el primer constituyente:
\textit{bici(cleta)senda} (= «senda para bicicletas»);
\textit{homo(sexual)fóbico} (= «el que padece fobia a los homosexuales»).}
etc.), \forma{auto-} representa un acortamiento de una palabra más extensa, el
compuesto \forma{automóvil}; de modo que la interpretación correcta es
\sig{pista de/para automóvil(es)}.

\begin{exe}
  \ex\label{ex:01-auto-estructura}
    \begin{xlist}
      \ex Palabra prefijada: \forma{retrato} > \forma{\textbf{auto}rretrato}
          (= «retrato de uno mismo»)
      \ex Palabra compuesta con acortamiento del primer constituyente:
          \forma{automóvil} + \forma{pista} > \forma{auto(móvil)pista}
          (= «pista para automóvil(es)»)
    \end{xlist}
\end{exe}

Ahora bien, ¿se pueden conectar estos planteos, focalizados y centrados en el
nivel de la palabra, con el uso real de la lengua y la interpretación en
contexto? Los siguientes ejemplos, tomados de la prensa, pueden servir para
mostrar cómo el hablante puede hacer uso eficaz de su conocimiento de la lengua
para generar y reconocer nuevas formas.

\begin{exe}
  \ex\label{ex:01-descomer}
    \begin{xlist}
      \ex\label{ex:01-descomer-a} «Contratar y despedir debería ser natural como
          \forma{comer y descomer}». El debate por una presunta estrategia de
          «flexibilización laboral» por parte del Gobierno se deslizó en las
          últimas horas en medio de polémicas declaraciones de funcionarios del
          área de Trabajo. El secretario de Empleo y ex CEO del grupo Techint,
          Miguel Ángel Ponte, opinó que incorporar o despedir personal debería
          ser para las empresas natural como «comer y descomer»
          (\url{https://www.perfil.com/noticias/politica/ponte-contratar-y-despedir-deberia-ser-natural-como-comer-y-descomer.phtml},
          09/01/17).
      \ex\label{ex:01-descomer-b} \textsc{Flexibilización laboral, pasantías y
          reducción de aportes patronales}. La teoría «comer y descomer». El ex
          Techint y actual secretario de Empleo, Miguel Ponte, ilustró con que
          la movilidad en el mercado laboral tiene que ser similar al proceso de
          ingesta y evacuación de comida
          (\url{https://www.pagina12.com.ar/tags/empleo},
          21/01/17).\pendiente{Fecha confusa en el original}
    \end{xlist}
\end{exe}

El ejemplo \exref{ex:01-descomer-a} incluye una palabra, \forma{descomer}, que
en su momento llamó la atención de la ciudadanía en relación con el planteo
respecto de la flexibilización laboral.\footnote{La flexibilización laboral o
desregulación del mercado de trabajo es un modelo que elimina regulaciones
para contratar y despedir empleados por parte de las empresas y organizaciones
privadas.} Este desconcierto se refleja en \exref{ex:01-descomer-b}, donde se
trata de descifrar el significado de la palabra a partir de su relación con su
opuesto reverso \parencite[18]{cruse2001} \forma{comer}. Si \forma{comer} es
\sig{ingerir alimentos}, el proceso inverso, como bien se aclara en
\exref{ex:01-descomer-b}, es su \forma{evacuación} o «expulsión de excrementos
del interior del cuerpo». Ahora bien, ¿si ya existe en la lengua una palabra
como \forma{evacuación} para dar cuenta de este proceso, por qué se elige
«crear» otra? Indudablemente en una cuestión sensible para la sociedad, como la
flexibilización laboral, el hablante busca atenuar la expresión evitando un
término tabú y creando otro, que en este contexto cuenta como un eufemismo.

\begin{recuadro}[Para tener en cuenta]
Aun elementos menores, del nivel subléxico, como los afijos, pueden hacer
aportes significativos a la comprensión del texto, por lo que, además de
considerar sus significados descontextualizados, es importante integrarlos al
sentido global del texto, el cual es producto de las relaciones y oposiciones
que cada constituyente entabla con el resto en un contexto determinado
\parencite{giammatteo2015}.
\end{recuadro}

\subsection{Del nivel léxico a las relaciones sintagmáticas}

\subsubsection{Las conexiones sintácticas entre las palabras}

Para el hablante común, la palabra presenta una «visibilidad» de la que carecen
otras unidades lingüísticas, como podrían ser el morfema o la oración, lo que le
confiere una mayor apariencia de realidad psicológica y prominencia cognitiva
en la conciencia \parencite{julien2007}. Además, como sostiene
\textcite[49]{jakobson1967}, los ítems léxicos son entidades de doble vía
porque «están ligados necesariamente con el código por una relación interna y
con el mensaje por una relación externa». Esta doble relación implica dos
posibles orientaciones:

\begin{itemize}
  \item \textit{hacia el interior del sistema lingüístico}, las palabras
        constituyen una interfaz \parencite{jackendoff1990} que conecta todos
        los componentes de la estructura lingüística ---fonológico,
        morfosintáctico, semántico---; y
  \item \textit{hacia el exterior}, las palabras se proyectan en los múltiples
        mensajes con los que nos comunicamos, relacionando así el sistema
        lingüístico con el conocimiento de mundo.
\end{itemize}

Dado que son los elementos léxicos los que se proyectan sintagmáticamente para
formar unidades mayores (v.~§\ref{sec:01-proyeccion}), también en este nivel
podemos valernos de comparaciones para identificar las reglas que rigen la
formación de sintagmas y oraciones. Así, \exref{ex:01-concordancia} y
\exref{ex:01-modo} muestran la importancia fundamental de las categorías
gramaticales que, aunque se manifiestan morfológicamente en el interior de la
palabra, su posición periférica en ella indica que su función va más allá de la
palabra, ya que consiste en establecer conexiones entre los constituyentes
sintagmáticos (v. capítulo~4: §9):

\begin{exe}
  \ex\label{ex:01-concordancia}
    \begin{xlist}
      \ex L-o-s extenuad-o-s viajer-o-s recorri-ero-n el parque.
      \ex[*]{El extenuad-o viajer-o-s recorriero-n el parque.}
      \ex[*]{L-o-s extenuad-o-s viajer-o-s recorri-ó el parque.}
    \end{xlist}
  \ex\label{ex:01-modo}
    \begin{xlist}
      \ex\label{ex:01-modo-a} Creo que vienen / \juicio{*}{vengan}.
      \ex\label{ex:01-modo-b} No creo que vengan / \juicio{*}{vienen}.
    \end{xlist}
\end{exe}

En \exref{ex:01-concordancia}, los rasgos \sig{masculino} y \sig{plural} del
sustantivo \forma{viajer-o-s} se repiten en todos los elementos que integran el
SN, como el artículo \forma{l-o-s} y el adjetivo \forma{extenuad-o-s};
asimismo, el rasgo \sig{plural} también se retoma en el verbo
\forma{recorriero-n}, como señal de su vinculación con el núcleo del sujeto
\forma{viajero-s}.\pendiente{El original dice «‘masculino’ y ‘singular’»} En
\exref{ex:01-modo} el modo verbal es una categoría exigida por la correlación
entre el verbo de la oración principal ---\forma{creo}--- y el de la
subordinada: cuando el verbo principal es afirmativo requiere modo indicativo
\forma{vienen} \exref{ex:01-modo-a} y rechaza el subjuntivo; y viceversa
respecto de la forma negativa del verbo, que exige subjuntivo \forma{vengan}
\exref{ex:01-modo-b}. Cualquier alteración de estas correlaciones convierte a
las estructuras resultantes en agramaticales.

\subsubsection{La interfaz sintaxis-semántica}

En relación con las oraciones, es necesario considerar que conforman
estructuras típicas, a las que denominamos \termino{esquemas oracionales}
transitivos, intransitivos, reflexivos, pasivos, impersonales, etc. Estas
configuraciones se encargan de transmitir, mediante los recursos
léxico-gramaticales de la lengua, las representaciones mentales de los
distintos sucesos y situaciones que experimentamos, deseamos o simplemente
imaginamos. Se generan en torno a los verbos, que son las palabras que refieren
a los sucesos que se describen en la oración. En consecuencia, son las
exigencias de los verbos las que determinan la configuración oracional
resultante (aunque, según veremos, un mismo verbo puede adoptar distintas
configuraciones).

Teniendo lo anterior como punto de partida, podemos plantear que toda oración
se organiza en torno a tres estructuras interrelacionadas:

\begin{itemize}
  \item \termino{la estructura eventiva}, que tiene en cuenta el tipo de verbo
        que constituye el núcleo de la oración: de estado, de acción, de
        proceso,\footnote{V. nota~\ref{fn:01-accion}.} etc.;
  \item \termino{la estructura argumental}, constituida por las funciones
        semánticas, conocidas como papeles temáticos, que asumen los argumentos
        que el verbo exige en función de su significado: agente, tema/paciente,
        experimentante, etc.;
  \item \termino{la estructura sintáctica}, que da cuenta de las funciones que
        cumplen los distintos constituyentes presentes en la
        oración.\footnote{La estructura sintáctica y su relación con la
        estructura argumental se trata en el capítulo~7 y la estructura
        eventiva, en el capítulo~8.}
\end{itemize}

Veamos solo algunos ejemplos sencillos para ilustrar estos planteos:

\begin{exe}
  \ex\label{ex:01-argumentos}
    \begin{xlist}
      \ex\label{ex:01-argumentos-a} Tres hombres robaron una tienda de productos electrodomésticos.
      \ex\label{ex:01-argumentos-b} Un intendente del conurbano entregó viviendas a los habitantes del municipio.
      \ex\label{ex:01-argumentos-c} Los cimientos cedieron y la casa se derrumbó.
    \end{xlist}
\end{exe}

Tanto en \exref{ex:01-argumentos-a} como en \exref{ex:01-argumentos-b} se trata
de verbos de \sig{acción} ---\forma{robar, entregar}---, pero mientras el
primero exige solo dos participantes, un agente que realiza la acción
---\forma{tres hombres}--- y un tema afectado por el suceso ---\forma{una tienda
de productos electrodomésticos}---, el segundo verbo, además del agente
---\forma{un intendente del conurbano}--- y del tema ---(las)
\forma{viviendas}---, requiere también un tercer participante en el evento, la
meta o término de la acción ---\forma{los habitantes del municipio}---. Pero
¿qué ocurre en \exref{ex:01-argumentos-c}? Los sucesos que se describen no se
presentan como resultado de la acción de ningún agente, sino que se trata de
procesos ---\forma{cedieron} y \forma{se derrumbó}--- que afectan a un solo
argumento cada uno respectivamente ---\forma{los cimientos} y \forma{la
casa}---, que cumplen la función de temas afectados por los eventos. Si bien no
podemos extendernos mayormente en estas consideraciones, que serán desarrolladas
en los capítulos correspondientes de este libro, los ejemplos presentados
muestran que los verbos difieren en sus exigencias argumentales y, lo que
resulta más interesante, es que estas diferencias no son aleatorias, sino que,
en gran medida, derivan de sus diferencias semánticas.

Asimismo, desde la perspectiva sintáctica, estos significados que conforman la
estructura argumental del verbo se reflejan en la organización de la estructura
oracional, estableciendo configuraciones canónicas o por defecto para los
distintos tipos de verbos. Así, los verbos de \sig{acción}, como \forma{robar}
en \exref{ex:01-argumentos-a}, típicamente involucran dos argumentos, agente y
tema, que funcionan, respectivamente, como sujeto y COD. \forma{Entregar}, en
cambio, forma parte del subgrupo de verbos de \sig{acción} que implican una
\sig{transferencia}, integrado también por verbos como \forma{dar, prestar,
regalar} o \forma{ayudar}, entre otros. Los verbos de \sig{transferencia} son
ditransitivos (o doblemente transitivos) y exigen, como vimos en
\exref{ex:01-argumentos-b}, tres argumentos: el agente y el tema funcionan
sintácticamente igual que con los transitivos, es decir, como sujeto y COD
respectivamente; pero el tercer argumento, la meta, cumple la función
sintáctica de COI. En el caso de los verbos de \exref{ex:01-argumentos-c}, que
manifiestan procesos que se presentan como no agentivos, el único argumento
exigido es el tema afectado, que en estos casos cumple la función sintáctica de
sujeto del verbo, como lo muestra en \exref{ex:01-argumentos-c} la
concordancia: plural en \forma{cedieron} porque el sujeto es \forma{los
cimientos}, y singular en \forma{se derrumbó}, cuyo sujeto es \forma{la casa}.

En función de los planteos que estamos presentando de modo ilustrativo y no
exhaustivo, ya que la mayoría se irá ampliando a lo largo de los capítulos que
siguen, resulta también importante resaltar que las distintas estructuras
oracionales poseen, además de la forma canónica, otras configuraciones
alternativas que le permiten al hablante mostrar distintas «visiones» o
perspectivas respecto de los sucesos que presenta. Dentro de estas
posibilidades de alternancia la más reconocida es la que se da entre la voz
activa y la pasiva, que en español puede manifestarse utilizando tanto la
pasiva perifrástica \exref{ex:01-pasivas-a}\footnote{La voz pasiva y demás
perífrasis verbales del español se desarrollan en el capítulo~11; y la pasiva
con \textit{se} y otras construcciones pronominales, en el capítulo~10.} como
la llamada pasiva con \forma{se} \exref{ex:01-pasivas-b}. Ambas pasivas
permiten promover al argumento tema/paciente\footnote{Ambos, tema y paciente,
representan al mismo argumento, pero se emplea la denominación de
\textit{tema} cuando el sustantivo que lo manifiesta posee el rasgo \sig{no
animado}; y la de \textit{paciente}, cuando es \sig{animado}.} a la posición
más prominente de la oración: la de sujeto, y mover al agente a una posición
marginal \exref{ex:01-pasivas-a} o incluso eliminarlo
\exref{ex:01-pasivas-b}.

\begin{exe}
  \ex\label{ex:01-pasivas}
    \begin{xlist}
      \ex\label{ex:01-pasivas-a} Una tienda de productos electrónicos fue robada (por tres hombres).
      \ex\label{ex:01-pasivas-b} Se entregaron viviendas a los habitantes del municipio
          \juicio{*}{por el intendente}.
    \end{xlist}
\end{exe}

Además de la alternancia activa/pasiva, la lengua presenta para los verbos
distintas formas de construcción correlativas de diferentes maneras de enfocar
los sucesos, tal como se muestra en los ejemplos de \exref{ex:01-perspectiva}:

\begin{exe}
  \ex\label{ex:01-perspectiva}
    \begin{xlist}
      \ex\label{ex:01-perspectiva-a} Un fuerte temporal de lluvia y viento azotó varios distritos de
          la provincia de Buenos Aires tras la intensa ola de calor que sufrió
          gran parte del país
          (\url{https://www.lanacion.com.ar/sociedad/temporal-en-chascomus-la-desesperacion-de-la-gente-que-fue-sorprendida-por-la-tormenta-nid},
          16/01/22).
      \ex\label{ex:01-perspectiva-b} Se derrumbó una casa en Rosario y buscan a un hombre de 52 años
          perdido entre los escombros
          (\url{https://www.perfil.com/noticias/actualidad/se-derrumbo-una-casa-en-rosario-y-buscan-a-un-hombre-de-52-anos-perdido-entre-los-escombros.phtml},
          26/07/22).
      \ex\label{ex:01-perspectiva-c} Las casas fabricadas con este material se construyen
          rápidamente y se necesita personal especializado para colocar
          tornillos y soldar las uniones
          (\url{https://mejorconsalud.as.com/casas-prefabricadas-debes-saber-comprar-una/}).
    \end{xlist}
\end{exe}

Si comparamos \exref{ex:01-perspectiva-a} con \exref{ex:01-argumentos-c}, en la
que los dos verbos refieren procesos que se presentan como no agentivos, el
primero de ellos, \exref{ex:01-perspectiva-a}, introduce una variante de este
papel temático, la causa ---\forma{un fuerte temporal de lluvia y viento}---,
que, en este caso, representa al fenómeno natural que provoca el hecho y que,
además, puede cumplir la función sintáctica de sujeto de la oración. En
\exref{ex:01-perspectiva-b}, el verbo transitivo \forma{derrumbar} se
intransitiviza mediante la pronominalización que lo convierte en
\forma{derrumbarse}, lo que habilita una lectura según la cual el suceso se
presenta como no provocado por ningún agente o causa, de modo que la oración
puede interpretarse \forma{La casa se derrumbó sola}. En
\exref{ex:01-perspectiva-c}, aunque \forma{construir} sea también un verbo de
\sig{acción} transitivo que supone un agente, se diferencia de
\exref{ex:01-perspectiva-b} porque hay en la oración ciertas marcas gramaticales
específicas, como el uso del presente y la presencia de un adverbio con valor
modal, que hacen que el agente se interprete como indeterminado y la lectura sea
de tipo genérico: \forma{Cualquiera puede construir rápidamente / en un abrir y
cerrar de ojos estas casas}.

Todas estas diferentes posibilidades de construcción, que configuran los
distintos esquemas oracionales estables de la lengua, los cuales pueden ser
reconocidos a partir de las marcas léxicas y gramaticales con que se
construyen, no pueden ser ignoradas cuando estudiamos una lengua. No se trata
solo de describir estructuras y de realizar minuciosos análisis de sus
componentes, sino de reflexionar acerca de la amplia gama de matices de
significación que estas estructuras permiten manifestar; se trata, en suma, de
que el conocimiento de estos esquemas sintácticos no se reduzca a un tema a
memorizar para la evaluación de la materia, sino de que se integre al bagaje de
estrategias del que podamos disponer como hablantes competentes de la lengua,
ya que este tipo de estructuras son las que a diario encontramos en los textos
con los que debemos interactuar.

\section{El sistema en uso: las pistas gramaticales en el texto}

En relación con el último eje planteado no vamos a hacer aquí un desarrollo
extenso, pero, al menos, nos interesa aportar un ejemplo, que nos permitirá
comprobar si el enfoque que estamos planteando, al considerar conjuntamente el
funcionamiento de los ítems léxicos y las estructuras morfosintácticas, resulta
adecuado no solo para aumentar el poder descriptivo y explicativo del análisis,
sino también para orientar la comprensión lectora y «convertir el saber
conceptual sobre el sistema de la lengua, en uso en géneros discursivos
concretos» \parencite{fontich2006}.

Para ilustrar estos conceptos, el texto elegido es el microrrelato de Jorge
Luis Borges \textit{El cautivo}, sobre el que no proponemos realizar un
análisis exhaustivo, sino simplemente poner de relieve cómo los recursos
léxico-gramaticales de la lengua intervienen en su construcción y, por tanto,
para su comprensión resulta útil prestar atención a su funcionamiento. Leamos
primero el texto:

\begin{textolectura}[El cautivo, de J. L. Borges]
En Junín o en Tapalqué refieren la historia. Un chico desapareció después de un
malón; se dijo que lo habían robado los indios. Sus padres lo buscaron
inútilmente; al cabo de los años, un soldado que venía de tierra adentro les
habló de un indio de ojos celestes que bien podía ser su hijo. Dieron al fin con
él (la crónica ha perdido las circunstancias y no quiero inventar lo que no sé)
y creyeron reconocerlo.

El hombre, trabajado por el desierto y por la vida bárbara, ya no sabía oír las
palabras de la lengua natal, pero se dejó conducir, indiferente y dócil, hasta
la casa. Ahí se detuvo, tal vez porque los otros se detuvieron. Miró la puerta,
como sin entenderla. De pronto bajó la cabeza, gritó, atravesó corriendo el
zaguán y los dos largos patios y se metió en la cocina. Sin vacilar, hundió el
brazo en la ennegrecida campana y sacó el cuchillito de mango de asta que había
escondido ahí, cuando chico. Los ojos le brillaron de alegría y los padres
lloraron porque habían encontrado al hijo.

Acaso a este recuerdo siguieron otros, pero el indio no podía vivir entre
paredes y un día fue a buscar su desierto. Yo querría saber qué sintió en aquel
instante de vértigo en que el pasado y el presente se confundieron; yo querría
saber si el hijo perdido renació y murió en aquel éxtasis o si alcanzó a
reconocer, siquiera como una criatura o un perro, los padres y la casa.
\end{textolectura}

Empecemos por considerar algunas cuestiones relativas al encuadre general del
texto. \textit{El cautivo} se presenta como un relato construido por un
narrador no omnisciente que, como es típico en Borges, simula no conocer todo
lo que concierne a la historia y la mente de sus personajes. Esta perspectiva
asumida por el narrador da ocasión a que el relato tenga dos finales sucesivos.
En el segundo párrafo, siguiendo lo que describe a partir de las señales
externas del comportamiento de los personajes, la historia del cautivo tiene un
primer final feliz ---\forma{Los ojos le brillaron de alegría y los padres
lloraron porque habían encontrado al hijo}---. Pero en el párrafo siguiente, el
narrador decide continuar la historia y propone un segundo final, que se
contrapone al anterior, lo que se explicita mediante el uso de la conjunción
coordinante adversativa inicial, encargada de señalar el cambio
---\forma{\textbf{pero} el indio no podía vivir entre paredes y un día fue a
buscar su desierto}---.

En cuanto a su organización global, el cuento presenta: 1) un \textbf{encuadre
de referencia}, que introduce lo que se cuenta en la región acerca de la
historia del cautivo; 2) el \textbf{núcleo narrativo}, que relata la historia
de la pérdida, recuperación y fuga final del protagonista; y 3) la
\textbf{evaluación}, que contiene los comentarios subjetivos del narrador
acerca de la historia. Ahora bien, un primer aspecto a destacar es que las
distinciones anteriores se correlacionan con cambios en la categoría gramatical
de persona. En el encuadre de referencia, el efecto de indeterminación se
transmite mediante formas propias de la impersonalidad, como el uso de la
3.ª~persona plural con sujeto indeterminado ---\forma{refieren (la
historia)}---, la 3.ª~singular pasiva (pasiva con \forma{se})
---\forma{se dijo}--- y la inclusión de un sujeto no personal ---\forma{la
crónica (ha perdido)}---. En el núcleo narrativo, en el que se relatan los
hechos fundamentales de la historia, se emplean formas de 3.ª~persona singular
o plural, pero referidas a sujetos personales y determinados, correspondientes
a los personajes de la historia ---el cautivo, sus padres, un soldado---. En la
evaluación final, la perspectiva del narrador se introduce mediante la
1.ª~persona, resaltada por la presencia repetida del pronombre \forma{yo}. La
tabla~\ref{tab:01-cautivo} resume lo expuesto:

\begin{table}[htbp]
  \centering
  \footnotesize
  \begin{tabularx}{\linewidth}{@{}>{\raggedright\arraybackslash}p{.24\linewidth}
                                  >{\raggedright\arraybackslash}X
                                  >{\raggedright\arraybackslash}p{.24\linewidth}@{}}
    \toprule
    \textbf{Encuadre de referencia} & \textbf{Núcleo narrativo} & \textbf{Evaluación} \\
    \textsc{lo que cuentan}          & \textsc{la historia del cautivo} & \textsc{comentarios del narrador} \\
    \textit{Uso de formas impersonales} & \textit{Uso de formas de 3.ª pers. sg. y pl. referidas a sujetos personales} & \textit{Uso de la 1.ª pers.} \\
    \midrule
    \textbf{Refieren} la historia (3.ª impersonal) \newline \textbf{Se dijo} (pasiva con \textit{se})
      & \textit{Desaparición}: Un chico desapareció. Lo habían robado los indios. Sus padres lo buscaron inútilmente.
      & \\
    \addlinespace
    \textbf{La crónica} ha perdido las circunstancias (sujeto no personal)
      & \textit{Encuentro/Regreso}: Un soldado les habló de un indio\ldots{} (Los padres) Dieron al fin con él. El hombre se dejó conducir.
      & No quiero inventar lo que no \textbf{sé} \\
    \addlinespace
      & \textit{Reencuentro}: Se detuvo, miró la puerta\ldots{} Gritó, bajó la cabeza, atravesó corriendo, se metió en la cocina. Sin vacilar, hundió el brazo, sacó el cuchillito. Los ojos le brillaron, los padres lloraron.
      & \\
    \addlinespace
      & \textit{Desencuentro/Huida}: El hombre no podía vivir entre paredes. Un día el indio se fue a buscar su desierto.
      & \textbf{Yo} querría saber qué sintió en aquel instante de vértigo\ldots; \textbf{yo} querría saber\ldots \\
    \bottomrule
  \end{tabularx}
  \caption[Correlación de las personas gramaticales]{Correlación de las
  personas gramaticales y las personas gramaticales\protect\pendiente{Epígrafe
  repetido en el original}}\label{tab:01-cautivo}
\end{table}

Uno de los aspectos más destacados del relato es que se presenta atravesado por
la incertidumbre. Ahora bien, ¿cómo apoyan los elementos léxico-gramaticales
esta sensación que transmite el texto? El relato se inicia con una coordinación
disyuntiva, que no tiene el típico valor de \sig{elección} de estas
construcciones, como el que se observa, por ejemplo, en \forma{¿Querés un
helado o una bebida fría? ¿Vamos al cine o al teatro?} La disyunción inicial
\forma{en Junín o en Tapalqué}, que sugiere que los hechos sucedieron «en algún
lugar de La Pampa», genera el efecto de presentar como poco seguro el lugar en
que transcurren los hechos, a la vez que remite intertextualmente al comienzo
del \textit{Quijote} («En un lugar de la Mancha, de cuyo nombre no quiero
acordarme»). Los elementos léxicos también colaboran en señalar falta de
certeza: así, por ejemplo, encontramos la combinación \forma{bien podía ser},
en la que el adverbio equivale a \forma{tal vez}; el uso del verbo
\forma{creer} ---\forma{creyeron reconocerlo}---, que le quita seguridad al
reconocimiento; adverbios de duda como \forma{tal vez} y \forma{acaso}; la
construcción aproximativa \forma{como sin entenderla}. En la parte
correspondiente a la evaluación, también abundan las marcas de incertidumbre:
el uso reiterado del condicional seguido del verbo \forma{saber}, que enfatiza
el desconocimiento por parte del narrador ---\forma{yo querría saber}---; la
interrogativa indirecta que manifiesta lo que ignora pero desea saber
---\forma{qué sintió en aquel instante de vértigo}---; y finalmente las dos
subordinadas condicionales coordinadas ---\forma{si el hijo perdido renació y
murió en aquel éxtasis o si alcanzó a reconocer, siquiera como una criatura o
un perro, los padres y la casa}---, que señalan enfáticamente y con un dejo de
ironía el deseo del narrador no omnisciente de conocer los resortes íntimos de
la conducta de su personaje. La segunda condicional plantea una disyuntiva
tácita: \forma{si alcanzó a reconocer o no}, y una expectativa de mínima:
\forma{siquiera}. En contraste con el clima de incertidumbre que enmarca la
historia, es de destacar que dentro del núcleo narrativo el narrador incluye,
contrastivamente, una afirmación enfática: \forma{sin vacilar}, que es la que
impulsa al lector a aceptar el primer final feliz, que luego será desplazado
por el relato de la huida del indio al desierto.

Por último, también se destaca el papel que en este microrrelato cumplen los
SSNN con que se va sucesivamente nombrando al protagonista, los cuales, poco a
poco, van revelando su verdadera identidad. \forma{El cautivo} es la
denominación inicial del título, que luego no aparece dentro del texto, pero
que alude a la dualidad interna que atraviesa al protagonista: cautivo de
\sig{dos civilizaciones} o \sig{modos de vida}, tal como los presenta el
narrador. Dentro del texto, el protagonista es presentado inicialmente
precedido del determinativo\footnote{Los determinativos son palabras
funcionales que tienen un significado fundamentalmente gramatical, que se
precisa en su relación con otra(s) palabra(s). Los determinativos débiles
---\textit{un, alguno, muchos, varios, otros}, etc.--- sirven para introducir
entidades nuevas en el discurso, mientras que los determinativos fuertes
---\textit{los, mis, estos, todos, ambos}, etc.--- permiten hacer mención a
entidades conocidas previamente \parencite{giammatteo2009}.} indefinido, como
\forma{un chico}, según lo esperado en la primera mención de una entidad nueva
en el discurso. Luego, cuando el soldado cree descubrirlo, el texto vuelve a
utilizar otro sustantivo precedido por un determinativo indefinido ---\forma{un
indio de ojos azules}---, y la reiteración del indefinido hace que el sintagma
funcione como una segunda presentación del protagonista, ahora según «su
nueva» naturaleza adquirida «en el desierto y la vida bárbara». En la tercera
mención, el protagonista ya aparece como corresponde a una entidad que ya fue
mencionada en el discurso, por lo que va precedida del determinante definido
---\forma{el hombre}---. Este sintagma contrasta doblemente por el carácter
definido y por el sustantivo empleado, con la mención inicial ---\forma{un
chico}---, señalando, además, el paso del tiempo y el crecimiento del niño. En
el momento del reconocimiento, el narrador elige nombrarlo como \forma{el
hijo}, para ponerlo en correlación con la anterior referencia a sus padres
---\forma{y los padres lloraron porque habían encontrado al hijo}---. Se
manifiesta así que la serie denominativa \forma{el cautivo / un chico / un
indio / el hombre} culmina en \forma{el hijo} que los padres reencuentran.
Pero, como ya hemos dicho, la historia continúa y, en esta segunda parte, el
narrador emplea otro sintagma definido ---\forma{el indio}---, con el que
revela la verdadera identidad del protagonista, aquello en lo que se ha
convertido a través del cautiverio. Por último, en su planteo dubitativo final,
el narrador se refiere nuevamente al cautivo como \forma{el hijo}, pero le
añade el calificativo \forma{perdido}, que contrasta con la mención en que
aparecía como \forma{el hijo que los padres habían recuperado}. Así, el
narrador culmina la caracterización del personaje y completa la serie
denominativa: \forma{el cautivo / un chico / un indio de ojos azules / el
hombre / el hijo} (recuperado) / \forma{el indio / el hijo perdido}.

El análisis del texto no se ha agotado con el enfoque presentado; no obstante,
la perspectiva léxico-gramatical desarrollada, que parte del conocimiento del
sistema, pero que también se interesa por el uso, no solo hace posible
sistematizar el funcionamiento de los diferentes recursos lingüísticos en el
nivel oracional, sino que también posibilita considerar su proyección textual.

\begin{recuadro}[Reflexión final]
El saber gramatical cumple un papel esencial: hace explícito el conocimiento
tácito del hablante, permitiéndole distanciarse del uso cotidiano para comenzar
a reflexionar sobre la naturaleza del lenguaje. Asimismo, constituye un
conocimiento fundamental para la comprensión de los sutiles matices
involucrados en los múltiples mensajes que damos y recibimos cotidianamente. De
este modo, el reconocimiento de los valores expresivos que conjugan los
elementos léxicos y gramaticales, sin duda, contribuye a captar la intención
del productor y a una interpretación más plena del texto. Este conocimiento
reflexivo es de gran importancia porque permite entender cómo funciona la
lengua y, en consecuencia, operar más eficazmente con ella en el uso cotidiano.
En esta perspectiva, el aprendizaje de la lengua materna puede dejar de ser un
procedimiento estéril y repetitivo para convertirse en un espacio para el
desarrollo de las habilidades cognitivo-lingüísticas necesarias para acompañar
la lectura y estimular la producción de los textos requeridos en el proceso de
enseñanza-aprendizaje.
\end{recuadro}

% --- Secciones finales -------------------------------------------------
\lecturasrecomendadas

\solucionario

\begin{solucion}{\ref{ej:01-facultad}}
\begin{enumerate}[label=\arabic*., start=2, leftmargin=2em, itemsep=.2\baselineskip]
  \item Se trata de una facultad universal porque todos los seres humanos la
        poseen.
  \item Es privativa de nuestra especie porque, al menos hasta donde sabemos,
        ninguna otra especie ha desarrollado un lenguaje con las
        características del humano.
  \item El lenguaje sirve para la comunicación, pero también cumple otros
        fines, como desarrollar el pensamiento abstracto, organizar la actividad
        mental y planificar nuestra actuación futura.
  \item La forma oral del lenguaje es natural para el hombre porque es la que
        todos los humanos sin distinción poseen.
  \item Desde esta perspectiva, la escritura puede considerarse una técnica
        desarrollada para conservar las palabras, la cual debe ser aprendida.
  \item Sin el lenguaje la vida social no sería posible porque es básicamente
        el medio por el que planificamos y organizamos todas nuestras
        actividades comunitarias.
  \item El lenguaje es parte de la dotación biológica de nuestra especie porque
        es parte de nuestra herencia genética y, por lo tanto, estamos
        naturalmente capacitados para desarrollarlo.
  \item Para el desarrollo del lenguaje es necesario interactuar con el medio
        porque necesitamos del estímulo que proviene de escucharlo en el medio
        ambiente (de la misma manera que una semilla necesita la tierra y el
        agua para germinar en una planta).
  \item Somos capaces de producir combinaciones de unidades que constituyen
        nuevas oraciones que jamás hemos oído porque conocemos las unidades y
        sabemos cómo combinarlas (o bien: porque tenemos internalizada la
        gramática de nuestra propia lengua).
\end{enumerate}
\end{solucion}

\begin{solucion}{\ref{ej:01-combinaciones}}
\begin{enumerate}[label=\arabic*., start=2, leftmargin=2em, itemsep=.2\baselineskip]
  \item El V \forma{vivir} exige un sujeto \sig{animado} y no es compatible con
        un sujeto como \forma{la lluvia}.
  \item El V \forma{tener}, además de sujeto, exige un COD que dé cuenta de lo
        que el sujeto posee. Si no se explicita, la O queda incompleta.
  \item La P exige como complemento un SN, no un SA como \forma{triste}.
  \item El V \forma{rezar} exige un COD compatible con su significado
        ---\forma{una plegaria, un padrenuestro}, etc.---.
  \item La P \forma{en} exige un complemento referido a lugar ---\forma{en
        Roma, en la iglesia}--- o bien con el V \forma{rezar} (que es un V de
        \sig{habla}) podría referirse a un idioma ---\forma{en latín, en
        chino}---.
\end{enumerate}
\end{solucion}

\begin{solucion}{\ref{ej:01-diccionariamente}}
La palabra \forma{diccionariamente} que usa el científico Diego Golombek está
formada por el N \forma{diccionario} y el sufijo \forma{-mente}. Resulta
aceptable porque tiene interpretación en el contexto del artículo: se refiere a
\sig{la manera en que aprendemos las palabras en los diccionarios, es decir, en
forma descontextualizada}.
\end{solucion}

\referenciascapitulo
